\documentclass[11pt]{article}

\usepackage[margin=1.1in]{geometry}
\usepackage{amsmath,amssymb,amsthm,mathtools}
\usepackage[hidelinks]{hyperref}

\newtheorem{theorem}{Theorem}[section]
\newtheorem{lemma}[theorem]{Lemma}
\newtheorem{proposition}[theorem]{Proposition}
\newtheorem{corollary}[theorem]{Corollary}
\newtheorem{conjecture}[theorem]{Conjecture}
\theoremstyle{remark}
\newtheorem{remark}[theorem]{Remark}

\newcommand{\N}{\mathbb N}
\newcommand{\Rcal}{\mathcal R}
\newcommand{\eps}{\varepsilon}
\newcommand{\1}{\mathbf 1}
\newcommand{\lowerdens}{\underline d}
\newcommand{\upperdens}{\overline d}
\DeclareMathOperator{\lcm}{lcm}
\DeclareMathOperator{\supp}{supp}


\author{}
\date{}

\begin{document}

\begin{abstract}
Let $\mathcal R$ be the set of positive integers representable as the sum of an
initial segment of the positive divisors of some integer, and, for
$n\in\mathcal R$, let $f(n)$ be the least such integer.  Erd\H{o}s asked whether
$f(n)=o(n)$, whether this holds for almost all $n$, and whether
$\limsup_{n\to\infty,\ n\in\mathcal R}f(n)/n=\infty$.  We use a
computer-assisted classification showing that
$\mathcal R=\mathbb N\setminus\{2,5\}$ and obtain quantitative estimates for
both tails of $f(n)/n$.

There is an absolute constant $c>0$ such that, uniformly in $X$ and for all
sufficiently small $\delta>0$,
\[
 \#\{n\leq X:n\in\mathcal R,\ f(n)\leq\delta n\}
 \ll \exp\!\left\{-\exp\bigl((1/\delta)^c\bigr)\right\}X.
\]
Conversely, for every fixed $\delta>0$, the inequality $f(n)\leq\delta n$
holds on a set of positive lower density.  At the opposite tail, as
$T\to\infty$,
\[
 \liminf_{X\to\infty}\frac1X
 \#\{n\leq X:n\in\mathcal R,\ f(n)>Tn\}
 \geq\left(\frac{e^{-\gamma}}2-o(1)\right)
 \frac{\log_3T}{(\log T)(\log_2T)}.
\]
Moreover, for every fixed $\eta>0$, at least $c_\eta X/\log_3X$
integers $X/2<n\leq X$ satisfy
\[
 \frac{f(n)}n>
 \exp\!\left[\left(\frac1{\sqrt2}-\eta\right)
 \sqrt{\log_3n\,\log_4n}\right]
\]
for all sufficiently large $X$.  Here $\log_j$ denotes the $j$-fold iterated
natural logarithm.
\end{abstract}


\section{Introduction and main results}\label{sec:introduction}

Throughout, $\N=\{1,2,3,\ldots\}$.  For $m\in\N$, let $\tau(m)$ denote the
number of positive divisors of $m$, and write those divisors as
\[
 1=d_1(m)<\cdots<d_{\tau(m)}(m)=m.
\]
An integer $N\in\N$ is \emph{represented} if
\[
 N=d_1(m)+\cdots+d_k(m)
\]
for some $m\in\N$ and some $1\leq k\leq\tau(m)$.  Let $\Rcal$ be the set
of represented integers.  For $N\in\Rcal$, define
\[
 f(N)=\min\bigl\{m\in\N:N=d_1(m)+\cdots+d_k(m)
       \text{ for some }1\leq k\leq\tau(m)\bigr\}.
\]
For a set $S\subseteq\N$, define its lower and upper asymptotic densities by
\[
 \lowerdens(S)=\liminf_{X\to\infty}\frac{\#(S\cap[1,X])}{X},\qquad
 \upperdens(S)=\limsup_{X\to\infty}\frac{\#(S\cap[1,X])}{X}.
\]

For $j\geq1$, define the iterated logarithms recursively by
\[
 \log_1 x=\log x,\qquad \log_{j+1}x=\log(\log_j x),
\]
whenever the right-hand side is defined.  We write $\sigma(n)=\sum_{d\mid n}d$ for the sum-of-divisors
function and $\psi(x)=\sum_{p^a\leq x}\log p$ for Chebyshev's function.
As usual, $\varphi$ denotes Euler's totient function, $\zeta$ the Riemann
zeta function, and $\gamma$ the Euler--Mascheroni constant.  For positive
integers $a_1,\ldots,a_r$, we write $[a_1,\ldots,a_r]$ for their least common
multiple.  For a prime
$p$, $v_p(n)$ denotes the exponent of $p$ in $n$, and $p^a\Vert n$ means
$v_p(n)=a$.  We write $\pi(x)$ for the number of primes not exceeding $x$.
We also write $P^+(n)$ and $P^-(n)$ for the largest and least prime factors
of $n$, with $P^+(1)=1$ and $P^-(1)=+\infty$.  Finally, $\1_S$ denotes the
indicator function of a set $S$.  For a map $H:\N\to\N$, we write
$H(\N)=\{H(n):n\in\N\}$ for its range.

For $j\geq0$, define
\[
 \sigma_j(m)=d_1(m)+\cdots+d_{\tau(m)-j}(m)
 \qquad(0\leq j<\tau(m)),
\]
and put $\sigma_j(m)=0$ for $j\geq\tau(m)$.  Thus every divisor-prefix
sum of $m$ is $\sigma_j(m)$ for a unique $j$.

The exact representability classification was established using a
computer-assisted argument, together with a verifier \cite{jif,jifRepo}:
\begin{equation}\label{eq:exact-representability}
 \Rcal=\N\setminus\{2,5\}.
\end{equation}
In particular, every integer $N\geq6$ is represented.

\begin{theorem}[Uniform lower-tail upper bound]
\label{thm:small-upper}
There is an absolute constant $c>0$ such that, for all sufficiently small
$\delta>0$ and all $X\geq1$,
\[
 \#\{N\leq X:N\in\Rcal,\ f(N)\leq\delta N\}
 \ll \exp\!\left\{-\exp\bigl((1/\delta)^c\bigr)\right\}X.
\]
Consequently, for every fixed $M>0$,
\[
 \#\{N\leq X:N\in\Rcal,\ f(N)\leq\delta N\}
 \ll_M \delta^M X
\]
uniformly in $X$.
\end{theorem}

\begin{theorem}[Positive lower tail]\label{thm:small-values}
For every $\delta>0$ there is a constant $c_\delta>0$ such that
\[
 \#\{N\leq X:N\in\Rcal,\ f(N)\leq\delta N\}\geq c_\delta X
\]
for all sufficiently large $X$.
\end{theorem}

\begin{theorem}[Almost-logarithmic fixed upper tail]
\label{thm:almost-log-tail}
As $T\to\infty$,
\begin{equation}\label{eq:almost-log-tail}
 \liminf_{X\to\infty}\frac1X
 \#\left\{N\leq X:
       N\in\Rcal,\ N\text{ squarefree},\ f(N)>TN
       \right\}
 \geq\left(\frac{e^{-\gamma}}2-o(1)\right)
       \frac{\log_3T}{(\log T)(\log_2T)}.
\end{equation}
In particular, the same lower bound holds without the squarefreeness
restriction.
\end{theorem}

\begin{theorem}[Subexponential squarefree growing levels]
\label{thm:subexp-growth}
For every fixed $\eta>0$ there exists $c_\eta>0$ such that, for all
sufficiently large $X$,
\begin{equation}\label{eq:subexp-growth}
\begin{aligned}
 &\#\biggl\{\frac X2<N\leq X:
       N\in\Rcal,\ N\text{ squarefree},\\
 &\hspace{38mm}
       \frac{f(N)}N>
       \exp\left[
         \left(\frac1{\sqrt2}-\eta\right)
         \sqrt{\log_3N\,\log_4N}
       \right]\biggr\}
 \geq c_\eta\frac X{\log_3X}.
\end{aligned}
\end{equation}
Consequently, for every fixed $s>0$ and $\eta>0$,
\begin{equation}\label{eq:positive-moment-growth}
 \frac1X\sum_{\substack{N\leq X\\N\in\Rcal}}
       \left(\frac{f(N)}N\right)^s
 \geq
 \exp\left[
       \left(\frac{s}{\sqrt2}-\eta\right)
       \sqrt{\log_3X\,\log_4X}
      \right]
\end{equation}
for all sufficiently large $X$.
\end{theorem}

Theorems~\ref{thm:small-upper} and~\ref{thm:small-values} give complementary
bounds for the lower tail: for each fixed $\delta>0$ it has positive lower
density, while its normalized counting function tends to zero, uniformly in
$X$, faster than any power of $\delta$ as $\delta\downarrow0$.  The proofs use,
respectively, a higher-moment estimate followed by Markov's inequality and a
collision estimate followed by Cauchy--Schwarz.

Theorem~\ref{thm:almost-log-tail} also settles the second question.  For
every fixed $T$ the set $\{N:f(N)>TN\}$ has positive lower density, so
$f(N)=o(N)$ fails on every set of density one.  It is still possible that
the density of $\{N:f(N)>TN\}$ tends to zero as $T\to\infty$.
Section~\ref{sec:tightness} shows that this is equivalent to a coverage
statement for bounded cofactors, and leaves it open.

The large-value results combine the large-cofactor moment estimate with a
rough-target argument.  Typically $\sigma(m)$ is divisible by every small
prime; if the represented value is rough, a small-cofactor representation is
therefore forced to contain a large divisor of a fixed least common multiple.
This leads to Theorems~\ref{thm:almost-log-tail} and~\ref{thm:subexp-growth}.

The exact classification \eqref{eq:exact-representability} and its verifier
first appeared in \cite{jif,jifRepo}.  The higher-moment estimate underlying
Theorem~\ref{thm:small-upper} first appeared in Kova\v{c}'s note \cite{Kovac},
extending the first-moment divisor-reflection argument in Tao's forum comment
\cite{TaoComment}.

Section~\ref{sec:tightness} studies the remaining density-one and limiting-law
questions.  It proves a first-moment result in every residue class, identifies
bounded-cofactor coverage exactly with tightness of the empirical measures,
and isolates the collision problem and the main open directions.

\section{Large values of the least representing integer}\label{sec:large-values}

\subsection{Preparatory estimates}

For a prime $q$ and a real number $y\geq1$, define
\[
 B_q(y)=\#\{n\leq y:q\nmid\sigma(n)\}.
\]
For integers $q\geq1$ and $a$ with $(a,q)=1$, let $\pi(x;q,a)$ denote the
number of primes $p\leq x$ satisfying $p\equiv a\pmod q$.  The following
estimate follows from a mean-value theorem for nonnegative multiplicative
functions and the Siegel--Walfisz theorem.

\begin{lemma}\label{lem:sigma-rate}
There are absolute constants $c_0,c_1>0$ such that, for all sufficiently
large $y$ and every prime
\[
 3\leq q\leq c_0\frac{\log_2 y}{\log_3 y},
\]
one has
\[
 B_q(y)\ll y\exp\!\left(-c_1\frac{\log_2 y}q\right).
\]
Also $B_2(y)\ll\sqrt y$.
\end{lemma}

\begin{proof}
For prime powers define the multiplicative function $h_q$ by
\[
 h_q(p^\nu)=
 \begin{cases}
  0,& p\equiv-1\pmod q\text{ and }\nu=1,\\
  1,& \text{otherwise}.
 \end{cases}
\]
If $p\equiv-1\pmod q$ and $v_p(n)=1$, then
$q\mid\sigma(p)\mid\sigma(n)$.  Hence
$\1_{q\nmid\sigma(n)}\leq h_q(n)$.  The hypotheses of the Hall--Tenenbaum mean-value estimate \cite[Theorem III.3.5]{Tenenbaum} hold with
absolute constants: the prime sum is bounded by Chebyshev's estimate, and
the contribution of prime powers of exponent at least two is convergent.
Thus,
\begin{align*}
 B_q(y)
 &\leq\sum_{n\leq y}h_q(n)\\
 &\ll y\prod_{p\leq y}
 \left(1-\frac1p+\sum_{\nu\geq1}\frac{h_q(p^\nu)}{p^\nu}\right)\\
 &\ll y\exp\!\left(
 -c\sum_{\substack{3\leq p\leq y\\p\equiv-1\ (q)}}\frac1p\right).
\end{align*}
Indeed, the factors at primes outside the indicated progression are
$1+O(p^{-2})$ and therefore have a bounded product.  For $p\geq3$ with
$p\equiv-1\pmod q$, the corresponding factor equals
$1-p^{-1}+[p(p-1)]^{-1}\leq\exp(-1/(2p))$.

Choose $z=\exp(Cq^2)$ with $C$ sufficiently large.  After decreasing
$c_0$ if necessary, the stated range for $q$ ensures that $z\leq y$ for
all sufficiently large $y$.  By the Siegel--Walfisz theorem \cite[Chapter II.8]{Tenenbaum}, uniformly for
$t\geq z$,
\[
 \pi(2t;q,-1)-\pi(t;q,-1)\gg\frac{t}{q\log t}.
\]
Summing over dyadic intervals from $z$ to $y$ therefore gives
\[
 \sum_{\substack{z<p\leq y\\p\equiv-1\ (q)}}\frac1p
 \gg \frac1q\log\frac{\log y}{\log z}
 =\frac1q\bigl(\log_2 y-2\log q+O(1)\bigr).
\]
In the stated range the last parenthesis is at least
$\tfrac12\log_2 y$ for large $y$, proving the first assertion.
Finally, $\sigma(n)$ is odd exactly when $n$ is a square or twice a square,
so $B_2(y)\ll\sqrt y$.
\end{proof}

For $t\geq2$ put $\mathfrak Q(t)=\prod_{p\leq t}p$, and for $u\geq1$ put
$\Lambda_u=\lcm(1,\ldots,u)$.  Mertens' theorem and Chebyshev's estimates give
\begin{equation}\label{eq:standard-prime-estimates}
 \prod_{p\leq t}\left(1-\frac1p\right)\asymp\frac1{\log t},
 \qquad \log\mathfrak Q(t)\ll t,
 \qquad \log\Lambda_u\leq C_\Lambda u
\end{equation}
for some absolute constant $C_\Lambda>0$; see, for example,
\cite[Chapter I.1]{Tenenbaum}.  The identity
$\log\Lambda_u=\psi(u)$ gives the last estimate.

For $e,d\geq1$, define
\[
 F_e(d):=\sum_{\substack{r\mid ed\\r\leq d}}r,
 \qquad
 g_e(m):=\sum_{\substack{r\mid m\\r\geq e}}\frac1r.
\]
If $N=d_1(m)+\cdots+d_k(m)$, let $d=d_k(m)$ be the final divisor
in the sum and put $e=m/d$.  Then $N=F_e(d)$; the integer $e$ will be called
the \emph{cofactor} of this representation.

\begin{lemma}[Bounded cofactors]\label{lem:bounded}
For all $e,d\geq1$,
\begin{equation}\label{eq:reflection}
 F_e(d)=ed\,g_e(ed).
\end{equation}
Moreover, if $m=ed$ and $N=F_e(d)$, then some prime
$p\leq\max(2,e\Lambda_e)$ satisfies
\begin{equation}\label{eq:prefix-congruence}
 N\equiv\sigma(m)\pmod p.
\end{equation}
Consequently, if $E,P,X\geq2$ and $E\Lambda_E\leq P$, then
\begin{equation}\label{eq:bounded-count}
 \#\{N\leq X:(N,\mathfrak Q(P))=1,\ N=F_e(d)\text{ for some }e\leq E\}
 \leq E\sum_{p\leq P}B_p(EX),
\end{equation}
where $\mathfrak Q(P)$ is the primorial defined above.
\end{lemma}

\begin{proof}
The map $r\mapsto ed/r$ exchanges the divisors $r\geq e$ of $ed$ with the
divisors at most $d$, proving \eqref{eq:reflection}.

For the congruence, let $\mathcal E=\{r:r\mid m,\ r<e\}$.  Complementary divisors give
\[
 \sigma(m)-N=m\sum_{r\in\mathcal E}\frac1r.
\]
If $\mathcal E$ is empty, take $p=2$.  If $\mathcal E=\{1\}$, any prime divisor of $e$ works.
Otherwise put $L=\lcm\{r:r\in\mathcal E\}$ and $C=\sum_{r\in\mathcal E}L/r$.  Then
$L\mid m$ and
\[
 \sigma(m)-N=\frac mL C.
\]
Since $1\in\mathcal E$ and $|\mathcal E|\geq2$, we have $C\geq2$; since
$L\leq\Lambda_e$ and $|\mathcal E|\leq e-1$, we also have
$C\leq(e-1)\Lambda_e<e\Lambda_e$.  Any prime divisor of $C$ therefore
proves \eqref{eq:prefix-congruence} with the stated bound.

Now suppose $e\leq E$, $N\leq X$, and $(N,\mathfrak Q(P))=1$.  Since the summand $d$ occurs in $F_e(d)=N$, we have $d\leq N$; hence the
integer $m=ed$ is at most $EX$.  Since the prime supplied above is at most $P$ and
$(N,\mathfrak Q(P))=1$, it does not divide $N$; the congruence then shows that it does
not divide $\sigma(m)$.  Counting the possible pairs $(e,m)$ and summing over
$p\leq P$ gives \eqref{eq:bounded-count}.
\end{proof}

The next estimate handles all large cofactors.  Its dependence on the moment
order will also be used when the threshold grows with $X$.

\begin{lemma}[Large cofactors]\label{lem:moment}
For every integer $k\geq2$ there is a constant $C_k$ satisfying
\begin{equation}\label{eq:Ck-growth}
 \log C_k\leq 2^{k+1}\log(2k+2)
\end{equation}
such that, for all $E,Z\geq1$,
\begin{equation}\label{eq:moment}
 \sum_{m\leq Z}\ \sum_{\substack{e\mid m\\e>E}}g_e(m)^k
 \leq C_kZE^{1-k}.
\end{equation}
Consequently, for all real $A,X\geq1$, the number of integers $N\leq X$
for which there are $e,d\in\N$ satisfying $N=F_e(d)$, $e>E$, and
$ed\leq AN$ is at most
\begin{equation}\label{eq:large-count}
 C_kA^{k+1}E^{1-k}X.
\end{equation}
\end{lemma}

\begin{proof}
Expanding the $k$th power and counting common multiples gives
\begin{align*}
 &\sum_{m\leq Z}\ \sum_{\substack{e\mid m\\e>E}}g_e(m)^k\\
 &\quad=\sum_{e>E}\ \sum_{r_1,\ldots,r_k\geq e}
 \frac1{r_1\cdots r_k}
 \left\lfloor\frac{Z}{\lcm(e,r_1,\ldots,r_k)}\right\rfloor\\
 &\quad\leq Z\sum_{e>E}\ \sum_{r_1,\ldots,r_k\geq e}
 \frac1{r_1\cdots r_k\,\lcm(e,r_1,\ldots,r_k)}.
\end{align*}
Set $x_0=e$, $x_i=r_i$, and $R=x_0\cdots x_k$.  Since
$r_1\cdots r_k\geq e^k$ and $e>E$,
\[
 \frac{R^{1/(k+1)}}{r_1\cdots r_k}
 \leq e^{1-k}\leq E^{1-k}.
\]
Consequently,
\[
 \frac1{r_1\cdots r_k\,\lcm(e,r_1,\ldots,r_k)}
 \leq E^{1-k}\frac{R^{-1/(k+1)}}{\lcm(x_0,\ldots,x_k)}.
\]
For a prime $p$ and $t\geq1$, let
$S_{p,t}=\{i:v_p(x_i)\geq t\}$, and for every nonempty
$S\subseteq\{0,\ldots,k\}$ put
\[
 y_S:=\prod_p p^{\#\{t\geq1:S_{p,t}=S\}}.
\]
Then
\[
 R=\prod_S y_S^{|S|},\qquad \lcm(x_0,\ldots,x_k)=\prod_S y_S.
\]
Dropping the nesting restrictions on the sets $S_{p,t}$ only enlarges the
sum, and hence
\begin{align*}
 \sum_{x_0,\ldots,x_k\geq1}
 \frac{R^{-1/(k+1)}}{\lcm(x_0,\ldots,x_k)}
 &\leq\prod_{\varnothing\ne S\subseteq\{0,\ldots,k\}}
 \sum_{y\geq1}y^{-1-|S|/(k+1)}\\
 &=\prod_{j=1}^{k+1}
 \zeta\!\left(1+\frac{j}{k+1}\right)^{\binom{k+1}{j}}
 =:C_k.
\end{align*}
This proves \eqref{eq:moment}.  Since
$\zeta(1+u)\leq1+u^{-1}$ for $u>0$, summing the binomial exponents gives
\eqref{eq:Ck-growth}.

For \eqref{eq:large-count}, put $m=ed$.  By \eqref{eq:reflection},
$N=mg_e(m)$.  Such a pair therefore satisfies
$m\leq AX$ and $g_e(m)=N/m\geq1/A$.  The number of distinct values $N$ is
at most the number of such pairs $(m,e)$, which by Markov's inequality and
\eqref{eq:moment} is at most
\[
 A^k\sum_{m\leq AX}\ \sum_{\substack{e\mid m\\e>E}}g_e(m)^k
 \leq C_kA^{k+1}E^{1-k}X.
\]
\end{proof}

\subsection{Sharper estimates for rough targets}

For readability in this subsection, for $y\geq2$ put
\[
 Q_y:=\mathfrak Q(y),\qquad
 \Delta_y:=\frac{\varphi(Q_y)}{Q_y},\qquad
 M_y:=\Lambda_{\lfloor y\rfloor}.
\]
Thus an integer is $y$-rough precisely when it is coprime to $Q_y$.

\begin{lemma}[Divisibility of divisor sums by a fixed modulus]
\label{lem:fixed-modulus-normality}
For every fixed positive integer $V$,
\[
 \#\{m\leq Y:V\nmid\sigma(m)\}=o_V(Y)
 \qquad(Y\to\infty).
\]
\end{lemma}

\begin{proof}
It suffices to prove the assertion when $V=q$ is a prime power.  For every
prime $r\equiv-1\pmod q$, the condition $r\Vert m$ implies
\[
 q\mid r+1=\sigma(r)\mid\sigma(m).
\]
Therefore, if $q\nmid\sigma(m)$, none of these primes divides $m$ exactly
once.  For any finite set $\mathcal P$ of such primes, the Chinese remainder
theorem gives
\[
 \limsup_{Y\to\infty}\frac1Y
       \#\{m\leq Y:q\nmid\sigma(m)\}
 \leq\prod_{r\in\mathcal P}
            \left(1-\frac1r+\frac1{r^2}\right).
\]
The reciprocal sum of the primes in the reduced residue class
$-1\pmod q$ diverges.  Letting $\mathcal P$ increase through this class makes
the displayed product tend to zero.  Factoring a general fixed $V$ into its
prime-power divisors and applying the union bound completes the proof.
\end{proof}

\begin{lemma}[Two divisor-tail estimates]
\label{lem:kernel-tails}
For $y,H\geq2$, put
\[
 S(y,H)=\sum_{\substack{L\mid M_y\\L>H}}\frac{\tau(L)}L.
\]
If $E\to\infty$ and
\[
 P=\left\lfloor
 E^{\,1+(2+(\log_3 E)^{-1/2})\log_2E/\log_3E}
 \right\rfloor,
\]
then
\begin{equation}\label{eq:fixed-kernel-tail}
 \log S(E,P/E)
 \leq-\frac{\log_2E}{\sqrt{\log_3E}}
      +O\left(\frac{\log_2E}{\log_3E}\right).
\end{equation}
If instead $J\to\infty$,
\[
 W=\sqrt{\frac{J\log J}{2}},
 \qquad F=\lfloor e^W\rfloor,
\]
then, for every fixed $C>0$,
\begin{equation}\label{eq:moving-kernel-tail}
 S\left(F,\frac{e^J}{CW}\right)
 \leq\exp\{-W+o(W)\}.
\end{equation}
\end{lemma}

\begin{proof}
For $0<\alpha\leq1/4$, Rankin's inequality and the Euler product for the
divisor function give
\begin{align}
 S(y,H)
 &\leq H^{-\alpha}
       \prod_{p\leq y}\sum_{a\geq0}
                \frac{a+1}{p^{a(1-\alpha)}}
  =H^{-\alpha}
       \prod_{p\leq y}(1-p^{-1+\alpha})^{-2}.
 \label{eq:rankin-kernel}
\end{align}
Mertens' estimate for $\sum_{p\leq y}1/p$, Chebyshev's bound, and partial
summation show that, when $z=\alpha\log y\to\infty$,
\begin{align}
 \sum_{p\leq y}p^{-1+\alpha}
 &=\log_2y+
   O\left(\int_0^z\frac{e^u-1}{u}\,du\right)
  =\log_2y+O\left(\frac{e^z}{z}\right).
 \label{eq:sharp-prime-sum}
\end{align}
The higher terms in the logarithm of the Euler product in
\eqref{eq:rankin-kernel} are $O(1)$ uniformly for $\alpha\leq1/4$.

For \eqref{eq:fixed-kernel-tail}, take
$\alpha=\log_3E/\log E$.  Then
\[
 \alpha\log(P/E)
 =\left(2+\frac1{\sqrt{\log_3E}}+o(1)\right)\log_2E,
\]
whereas \eqref{eq:sharp-prime-sum} gives
\[
 2\sum_{p\leq E}p^{-1+\alpha}
 =2\log_2E+O\left(\frac{\log_2E}{\log_3E}\right).
\]
Substitution in \eqref{eq:rankin-kernel} proves the first assertion.

For \eqref{eq:moving-kernel-tail}, take $\alpha=W/J$.  Since
\[
 \alpha\log F=\frac12\log J+o(1),
 \qquad
 \alpha\log\left(\frac{e^J}{CW}\right)=W+o(W),
\]
equation \eqref{eq:sharp-prime-sum} gives
\[
 \sum_{p\leq F}p^{-1+\alpha}
 \ll\log W+\frac{\sqrt J}{\log J}=o(W).
\]
The second assertion follows once more from \eqref{eq:rankin-kernel}.
\end{proof}

\begin{proof}[Proof of Theorem~\ref{thm:almost-log-tail}]
Fix an integer $j\geq3$ and put
\begin{align*}
 b_j&=\frac{j+1}{j-1},\\
 E&=\left\lceil T^{b_j}(\log T)^{3/(j-1)}\right\rceil,\\
 P&=\left\lfloor
 E^{\,1+(2+(\log_3E)^{-1/2})\log_2E/\log_3E}
 \right\rfloor.
\end{align*}
All parameters in this proof are fixed while $X\to\infty$.  For
$2\leq e\leq E$ and $L\mid M_E$, define
\[
 C_e(L)=\sum_{\substack{r\mid L\\r<e}}\frac Lr,
 \qquad
 k_e(L)=\frac{C_e(L)}{\operatorname{sm}_P(C_e(L))},
\]
where $\operatorname{sm}_P(v)$ is the largest divisor of $v$ all of whose
prime factors are at most $P$.  Set
\[
 K=\lcm_{\substack{2\leq e\leq E\\L\mid M_E}} k_e(L),
 \qquad V=KQ_P.
\]
Lemma~\ref{lem:fixed-modulus-normality} shows that the sources $m\leq EX$
with $V\nmid\sigma(m)$ number $o_{E,P}(X)$.  Since each source has at most
$E$ possible cofactors not exceeding $E$, their entire bounded-cofactor
image has cardinality $o_{E,P}(X)$.

Let $N\leq X$ be $P$-rough and suppose that
$N=F_e(m/e)$ with $e\leq E$ and $V\mid\sigma(m)$.  Since the summand $m/e$
occurs in $F_e(m/e)$, one has $m\leq eN\leq EX$.  The case $e=1$ is
impossible, since it would give $N=\sigma(m)$ although $Q_P\mid\sigma(m)$.
For $e\geq2$, put
\[
 L=\lcm\{r:r\mid m,\ r<e\},
 \qquad C=C_e(L),
 \qquad R=\sigma(m)-N.
\]
Reflection of the divisor list gives the exact identity
\begin{equation}\label{eq:complementary-kernel}
 R=\frac mL C.
\end{equation}
Since $Q_P\mid\sigma(m)$ and $(N,Q_P)=1$, the integer $R$ is $P$-rough.
Consequently both $u=m/L$ and $C$ are $P$-rough.  For every prime $p\mid e$,
one has $p\leq E<P$ and therefore
$v_p(L)=v_p(m)\geq v_p(e)$, because $m=Lu$ and $p\nmid u$.  Thus $e\mid L$.
Moreover $L\mid M_E$ and
\[
 C\geq L\geq e\geq2,
 \qquad C\leq EL.
\]
Thus $C>P$, $L>P/E$, and $k_e(L)=C$.  By the definition of $K$, one has
$C\mid K\mid\sigma(m)$, while \eqref{eq:complementary-kernel} gives
$C\mid R$.  Therefore $C\mid N$.

For fixed $e,L$, the $P$-rough multiples of $C$ have natural density
$\Delta_P/C$.  Since $C\geq L$ and $e\mid L$, summing over the finitely many
admissible pairs gives total upper density at most
\[
 \Delta_P\sum_{\substack{L\mid M_E\\L>P/E}}\frac{\tau(L)}L
 =\Delta_P S(E,P/E)=o(\Delta_P)
 \qquad(T\to\infty).
\]
By Lemma~\ref{lem:moment}, the targets admitting a representation with
$e>E$ and $m\leq TN$ have upper density at most
\[
 C_jT^{j+1}E^{1-j}
 \leq\frac{C_j}{(\log T)^3}=o(\Delta_P).
\]
Finally, a $P$-rough target which is not squarefree is divisible by $p^2$ for
some $p>P$, so these targets have upper density at most
\[
 \sum_{p>P}\frac1{p^2}\ll\frac1P=o(\Delta_P).
\]
By \eqref{eq:exact-representability}, at most two further targets are
unrepresented.  Subtracting these exceptional sets from the $P$-rough
integers gives
\[
 \liminf_{X\to\infty}\frac1X
 \#\{N\leq X:N\in\Rcal,\ N\text{ squarefree},\ f(N)>TN\}
 \geq(1-o(1))\Delta_P.
\]
Indeed,
\[
 \log E=(b_j+o(1))\log T,\qquad
 \log_2E=\log_2T+o(\log_2T),\qquad
 \log_3E=\log_3T+o(1),
\]
and hence
\[
 \log P=(2b_j+o(1))
          \frac{(\log T)(\log_2T)}{\log_3T}.
\]
Mertens' theorem therefore gives
\[
 \Delta_P
 =\left(\frac{e^{-\gamma}}{2b_j}+o(1)\right)
       \frac{\log_3T}{(\log T)(\log_2T)}.
\]
First let $T\to\infty$ for each fixed $j$, and then let $j\to\infty$.
Since $b_j\to1$, this proves \eqref{eq:almost-log-tail}.
\end{proof}

\begin{proof}[Proof of Theorem~\ref{thm:subexp-growth}]
Fix $0<\varepsilon<1$ sufficiently small in terms of $\eta$.  Choose a fixed
$\rho>0$ sufficiently small, depending only on the absolute constants in
Lemma~\ref{lem:sigma-rate}, and put
\[
 P=\left\lfloor\rho\frac{\log_2X}{\log_3X}\right\rfloor,
 \qquad J=\log P,
 \qquad W=\sqrt{\frac{J\log J}{2}},
 \qquad F=\lfloor e^W\rfloor,
 \qquad t=e^{(1-\varepsilon)W}.
\]
In particular $F<P$ for all sufficiently large $X$.  Complete periods modulo
$Q_P$ and Mertens' theorem give
\begin{equation}\label{eq:sharp-rough-target-count}
 \#\{X/2<N\leq X:(N,Q_P)=1\}
 =\frac X2\Delta_P+O(Q_P)
 \asymp\frac X{\log_3X},
\end{equation}
because $Q_P=\exp(O(P))=X^{o(1)}$ and $\Delta_P\asymp1/J$.  The
nonsquarefree members of this interval number at most
\[
 X\sum_{p>P}\frac1{p^2}\ll\frac XP=o(X\Delta_P).
\]

Write $Y=tX$.  Lemma~\ref{lem:sigma-rate} and the union bound show that the
number of sources $m\leq Y$ for which some prime $p\leq P$ fails to divide
$\sigma(m)$ is at most
\begin{equation}\label{eq:sharp-bad-source-count}
 O(\sqrt Y)+O\left(
      Y\pi(P)\exp\left[-c_1\frac{\log_2Y}{P}\right]
    \right).
\end{equation}
The first term treats $p=2$.  The chosen range of $P$ is permitted by
Lemma~\ref{lem:sigma-rate} after decreasing $\rho$ if necessary.  Since
\[
 Ft=e^{(2-\varepsilon)W}=(\log_2X)^{o(1)},
 \qquad
 \exp\left[-c_1\frac{\log_2Y}{P}\right]
      =(\log_2X)^{-c_1/\rho+o(1)},
\]
we may also ensure $c_1/\rho>3$.  Multiplying
\eqref{eq:sharp-bad-source-count} by the at most $F$ possible cofactors per
source gives $o(X\Delta_P)$ target--cofactor pairs.

Consider a remaining $P$-rough target $N$ admitting a witness
$N=F_e(m/e)$ with $m\leq tN$ and $e\leq F$.  Again $e=1$ is impossible.
For $e\geq2$, define $L,C,R$ as in \eqref{eq:complementary-kernel}.  Since
every prime $p\leq P$ divides $\sigma(m)$ but not $N$, the integer
$R=(m/L)C$ is $P$-rough.  Writing $u=m/L$, both $u$ and $C$ are therefore
$P$-rough.  Since $e\leq F<P$ and $e\mid m=Lu$, every prime-power divisor
of $e$ already divides $L$.  Consequently
\[
 m=Lu,\qquad (u,Q_P)=1,\qquad (C,Q_P)=1,
 \qquad e\mid L\mid M_F.
\]
As before $C\geq L\geq2$, so $C>P$.  On the other hand,
\[
 C\leq\sigma(L)
 \leq L\prod_{p\leq F}\left(1-\frac1p\right)^{-1}
 \ll L\log F.
\]
There is therefore an absolute constant $C_0>0$ such that
\begin{equation}\label{eq:moving-kernel-lower-bound}
 L>\frac P{C_0W}.
\end{equation}

Uniformly for $L\mid M_F$,
\[
 \log L\leq\log M_F\ll F=o(\log X),
 \qquad \log Q_P\ll P=o(\log X).
\]
It follows that $tX/L\geq Q_P$ for all sufficiently large $X$.  Complete
periods modulo $Q_P$ consequently give
\[
 \#\{u\leq tX/L:(u,Q_P)=1\}
 \leq\frac{2tX\Delta_P}{L}.
\]
For each fixed $L$, the admissible cofactors satisfy $e\mid L$, so there are
at most $\tau(L)$ possibilities.  Hence \eqref{eq:moving-kernel-lower-bound}
and Lemma~\ref{lem:kernel-tails} bound the entire remaining low-cofactor
witness count by
\begin{align*}
 2tX\Delta_P
     \sum_{\substack{L\mid M_F\\L>P/(C_0W)}}\frac{\tau(L)}L
 &\leq X\Delta_P\exp\{-\varepsilon W+o(W)\}\\
 &=o(X\Delta_P).
\end{align*}

Choose a fixed integer $k\geq2$ so large that $(k+1)\varepsilon>2$.
Lemma~\ref{lem:moment} bounds the targets with a witness $m\leq tN$ and
cofactor $e>F$ by
\[
 C_kt^{k+1}F^{1-k}X.
\]
After division by $X\Delta_P$, this is at most
\[
 O_k\left(J\exp\{[2-(k+1)\varepsilon+o(1)]W\}\right)=o(1).
\]
Thus a positive proportion of the squarefree targets in
\eqref{eq:sharp-rough-target-count} has no representation with $m\leq tN$.
By \eqref{eq:exact-representability}, all but at most two of these targets
are represented.  Finally,
\[
 \log t
 =\left(\frac{1-\varepsilon}{\sqrt2}+o(1)\right)
          \sqrt{\log_3X\,\log_4X}.
\]
Choose $\varepsilon$ sufficiently small in terms of $\eta$.  For
$X/2<N\leq X$, replacing $X$ by $N$ in these iterated logarithms changes
the right-hand side by $o(1)$ relative to its main term.  This proves
\eqref{eq:subexp-growth}.

For \eqref{eq:positive-moment-growth}, apply the first assertion with a
sufficiently small auxiliary value of $\eta$, raise the resulting ratios to
the power $s$, and note that
\[
 \log\log_3X=\log_4X
 =o\left(\sqrt{\log_3X\,\log_4X}\right).
\]
The factor $c_\eta/\log_3X$ in the proportion of surviving targets can
therefore be absorbed into the prescribed exponent loss $\eta$.
\end{proof}

\section{Small values of the least representing integer}\label{sec:small-values}

\subsection{Uniform upper bounds for the lower tail}

The following higher-moment estimate is due to Kova\v{c} \cite{Kovac},
extending the first-moment divisor-reflection argument in Tao's forum comment
\cite{TaoComment}.

\begin{lemma}[Kova\v{c}'s moment estimate]\label{lem:kovac-moment}
There is an absolute constant $C>0$ such that, for every integer $q\geq3$
and every real $x\geq1$,
\[
 \sum_{m\leq x}\sum_{j\geq0}
 \left(\frac{\sigma_j(m)}m\right)^q
 \ll \exp(Cq\log\log q)x.
\]
The implicit constant is absolute.
\end{lemma}

\begin{proof}
Let
\[
 1=r_1<r_2<\cdots<r_{\tau(m)}=m
\]
be the divisors of $m$.  Divisor reflection gives
\[
 \frac{\sigma_j(m)}m=\sum_{i>j}\frac1{r_i}.
\]
After expanding the $q$th power and summing over $j$, a fixed tuple
$a_1,\ldots,a_q$ of divisors is counted once for each divisor threshold not
exceeding $\min_i a_i$.  Hence
\[
 \sum_{j\geq0}\left(\frac{\sigma_j(m)}m\right)^q
 =
 \sum_{\substack{r,a_1,\ldots,a_q\mid m\\
                   r\leq\min(a_1,\ldots,a_q)}}
 \frac1{a_1\cdots a_q}.
\]
Summing over $m\leq x$ and bounding the number of multiples of
$[r,a_1,\ldots,a_q]$ by $x/[r,a_1,\ldots,a_q]$, it is enough to prove
\begin{equation}\label{eq:kovac-S}
 S:=\sum_{\substack{r,a_1,\ldots,a_q\geq1\\r\leq\min_i a_i}}
 \frac1{a_1\cdots a_q[r,a_1,\ldots,a_q]}
 \ll\exp(Cq\log\log q).
\end{equation}
For each $i$, put
\[
 b_i=\frac r{(r,a_i)},\qquad c_i=\frac{a_i}{(r,a_i)}.
\]
Then $b_i\mid r$, $(b_i,c_i)=1$, $a_i=rc_i/b_i$, and
\[
 [r,a_1,\ldots,a_q]=r[c_1,\ldots,c_q].
\]
The condition $r\leq a_i$ becomes $b_i\leq c_i$.  Dropping the coprimality
conditions and taking $\beta=(q-1)/q$, the inequality
\[
 \frac1{c_i}\leq\frac1{b_i^{1-\beta}c_i^\beta}
 \qquad(c_i\geq b_i)
\]
applies term by term.  Replacing each divisor $b_i\mid r$ by $r/d_i$
and using $\beta q=q-1$ leaves the power $r^{-2}$.  Thus
$S\leq S'S''$, where
\[
 S'=\sum_{r\geq1}\frac1{r^2}
 \left(\sum_{d\mid r}\frac1{d^{(q-1)/q}}\right)^q
\]
and
\[
 S''=\sum_{c_1,\ldots,c_q\geq1}
 \frac1{c_1^{(q-1)/q}\cdots c_q^{(q-1)/q}[c_1,\ldots,c_q]}.
\]
Both sums have Euler products.

For $S'$, write $S'=\prod_pL_p$, where
\[
 L_p=\sum_{\ell\geq0}p^{-2\ell}
 \left(\sum_{0\leq\gamma\leq\ell}
 p^{-((q-1)/q)\gamma}\right)^q.
\]
For $p\leq q$,
\[
 L_p\leq(1-p^{-2})^{-1}(1-p^{-(q-1)/q})^{-q}.
\]
Since
\[
 \sum_{p\leq q}p^{-(q-1)/q}
 \leq q^{1/q}\sum_{p\leq q}\frac1p
 \ll\log\log q,
\]
one obtains
\[
 \prod_{p\leq q}L_p\ll\exp(Cq\log\log q).
\]
For $p>q$, one has $qp^{-(q-1)/q}\ll1$, and the terms with $\ell\geq1$
give $L_p=1+O(p^{-2})$.  Thus
\begin{equation}\label{eq:kovac-Sprime}
 S'\ll\exp(Cq\log\log q).
\end{equation}

For $S''$, write $S''=\prod_pM_p$, where
\[
 M_p=\sum_{\gamma_1,\ldots,\gamma_q\geq0}
 p^{-((q-1)/q)(\gamma_1+\cdots+\gamma_q)-\max_i\gamma_i}.
\]
Put
\[
 A_h=\sum_{0\leq\gamma\leq h}p^{-((q-1)/q)\gamma}.
\]
Grouping by $h=\max_i\gamma_i$ gives
\[
 M_p=\sum_{h\geq0}p^{-h}(A_h^q-A_{h-1}^q),\qquad A_{-1}=0.
\]
For $p\leq q$,
\[
 M_p\leq(1-p^{-1})^{-1}(1-p^{-(q-1)/q})^{-q},
\]
so the product over these primes is $\ll\exp(Cq\log\log q)$.  For $p>q$,
the mean-value theorem and $qp^{-(q-1)/q}\ll1$ give
\[
 A_h^q-A_{h-1}^q\ll q p^{-h(q-1)/q}.
\]
Consequently,
\[
 M_p-1\ll q p^{-2+1/q},
\qquad
 \sum_{p>q}(M_p-1)\ll q\sum_{n>q}n^{-2+1/q}\ll1.
\]
Hence
\begin{equation}\label{eq:kovac-Ssecond}
 S''\ll\exp(Cq\log\log q).
\end{equation}
Combining \eqref{eq:kovac-Sprime} and \eqref{eq:kovac-Ssecond} proves
\eqref{eq:kovac-S} and the lemma.
\end{proof}

\begin{proof}[Proof of Theorem~\ref{thm:small-upper}]
Let
\[
 \mathcal E_\delta(X)=
 \{N\leq X:N\in\Rcal,\ f(N)\leq\delta N\}.
\]
If $\delta X<1$, then $\mathcal E_\delta(X)$ is empty, so assume
$\delta X\geq1$.  If $N\in\mathcal E_\delta(X)$ and $m=f(N)$, then
$m\leq\delta X$ and $N=\sigma_j(m)$ for some $j\geq0$ with
$\sigma_j(m)/m\geq1/\delta$.  Therefore, for every integer $q\geq3$,
\begin{align*}
 \#\mathcal E_\delta(X)
 &\leq\delta^q\sum_{m\leq\delta X}\sum_{j\geq0}
       \left(\frac{\sigma_j(m)}m\right)^q\\
 &\ll\delta^{q+1}\exp(Cq\log\log q)X.
\end{align*}
Choose a sufficiently small absolute $c>0$ and put
\[
 q=\left\lfloor\exp((1/\delta)^c)\right\rfloor.
\]
For all sufficiently small $\delta$, one has
$C\log\log q\leq\tfrac12\log(1/\delta)$, and hence
\[
 \delta^{q+1}\exp(Cq\log\log q)
 \leq\exp\!\left(-\frac q2\log\frac1\delta\right)
 \leq\exp\!\left\{-\exp((1/\delta)^{c/2})\right\}.
\]
Renaming $c/2$ as $c$ proves the first assertion.  The fixed-power bound
follows because this double-exponential expression is $O_M(\delta^M)$ for
every fixed $M>0$.
\end{proof}

\subsection{Positive-density lower bounds}

We now prove Theorem~\ref{thm:small-values}.

For $m\in\N$, define
\[
 s(m)=\sigma(m)-m.
\]
If $m\geq2$, then $s(m)$ is the sum of all proper divisors of $m$, so
$s(m)\in\Rcal$ and
\begin{equation}\label{eq:sv-basic}
  f(s(m))\leq m \qquad(m\geq2).
\end{equation}


By \eqref{eq:sv-basic}, it suffices to construct $\mathcal A(X)\subset[1,X]$
such that $s(m)>m/\delta$ on $\mathcal A(X)$ and
$\#s(\mathcal A(X))\gg_\delta X$.

It is enough to prove the theorem for $0<\delta\leq1$.  Put
$\lambda=1/\delta$, and choose a fixed squarefree integer $D$ such that
\begin{equation}\label{eq:sv-D}
  \frac{\sigma(D)}D>1+\lambda.
\end{equation}
Such a $D$ exists because $\prod_{p\leq z}(1+1/p)$ is unbounded.

\subsection{A family of uniformly abundant integers}

For large $X$, let $\mathcal A_0(X)$ consist of the integers
\begin{equation}\label{eq:sv-factorization}
  M=pqrk
\end{equation}
for which $p,q,r$ are prime,
\begin{align*}
  &X^{1/120}<k=Dj\leq X^{1/60},
       \qquad \frac{\sigma(j)}j\leq4\zeta(2),\\
  &X^{1/15}<r\leq X^{1/12},\\
  &X^{7/20}<q\leq X^{11/30},\\
  &\frac{X}{2qrk}<p\leq\frac{X}{qrk}.
\end{align*}
Set $\ell=rk$ and $m=q\ell$.

\begin{lemma}\label{lem:sv-family}
There is a constant $C_\delta>0$ such that, for all sufficiently large
$X$,
\[
  \#\mathcal A_0(X)\gg_\delta X,
\]
every $M\in\mathcal A_0(X)$ has a unique factorization
\eqref{eq:sv-factorization}, and
\begin{equation}\label{eq:sv-two-sided}
  \lambda<\frac{s(t)}t<C_\delta
  \qquad(t\in\{k,\ell,m,M\}).
\end{equation}
Moreover, $q>m^{7/9}$.
\end{lemma}

\begin{proof}
For $J\geq1$,
\[
 \sum_{j\leq J}\frac{\sigma(j)}{j^2}
 =\sum_{d\leq J}\frac1{d^2}\sum_{a\leq J/d}\frac1a
 \leq\zeta(2)\sum_{a\leq J}\frac1a.
\]
Consequently, with $J_1=X^{1/120}/D$ and $J_2=X^{1/60}/D$,
\begin{align*}
 \sum_{\substack{J_1<j\leq J_2\\ \sigma(j)/j\leq4\zeta(2)}}\frac1j
 &\geq \sum_{J_1<j\leq J_2}\frac1j
   -\frac1{4\zeta(2)}\sum_{j\leq J_2}\frac{\sigma(j)}{j^2}\\
 &\geq \left(\frac1{240}+o(1)\right)\log X.
\end{align*}
For fixed $q,r,k$, the interval for $p$ is $(T/2,T]$, where
$T=X/(qrk)\geq X^{8/15}$.  The prime number theorem, followed by Mertens'
theorem for primes, therefore gives
\[
 \#\mathcal A_0(X)
 \gg \frac{X}{\log X}
 \sum_k\frac1k
 \sum_{X^{1/15}<r\leq X^{1/12}}\frac1r
 \sum_{X^{7/20}<q\leq X^{11/30}}\frac1q
 \gg_\delta X.
\]
The factorization is unique because $p>X^{1/2}$, the prime $q$ lies above
$X^{7/20}$, the prime $r$ lies between $X^{1/15}$ and $X^{1/12}$, and
every prime factor of $k$ is at most $X^{1/60}$.

The function $n\mapsto\sigma(n)/n$ is nondecreasing under divisibility.
Since $D\mid t$ for $t\in\{k,\ell,m,M\}$, \eqref{eq:sv-D} therefore
gives the lower bound in \eqref{eq:sv-two-sided}.  The inequality
$\sigma(ab)/(ab)\leq(\sigma(a)/a)(\sigma(b)/b)$ also gives
\[
 \frac{\sigma(k)}k\leq
 \frac{\sigma(D)}D\frac{\sigma(j)}j
 \leq4\zeta(2)\frac{\sigma(D)}D,
\]
and adjoining each of $p,q,r$ multiplies $\sigma(t)/t$ by less than $2$.
This gives the upper bound in \eqref{eq:sv-two-sided}.  Finally,
$\ell\leq X^{1/10}$ and $q^2>X^{7/10}\geq\ell^7$, whence
$q^9>q^7\ell^7=m^7$.
\end{proof}

For $u$ sufficiently large, put
\[
  y(u)=\frac{\log_2u}{\log_3u},
  \qquad Y=y(X^{1/120}).
\]
An integer is called $Y$-smooth if all of its prime factors are at most
$Y$.

\begin{lemma}[Luca--Pomerance regularity inputs]\label{lem:LP-inputs}
Put $y(n)=\log_2n/\log_3n$.  Outside a set of integers of asymptotic density
zero, the following hold:
\begin{enumerate}
\item $v_p(\sigma(n))>v_p(n)$ for every prime $p\leq y(n)$;
\item $P^+(\gcd(n,\sigma(n)))\leq y(n)$;
\item $\sigma(n)/\gcd(n,\sigma(n))$ is divisible by every prime
      $p\leq y(n)$;
\item every prime factor of $s(n)/\gcd(n,\sigma(n))$ exceeds $y(n)$.
\end{enumerate}
If in addition $P^+(n)>n^{7/9}$, then, outside another density-zero set,
\[
 y(n)<\pi\leq n^{10/27},\quad \pi\text{ prime}
 \quad\Longrightarrow\quad \pi^2\nmid s(n).
\]
Finally, outside a set of asymptotic density zero,
\[
 \sum_{\substack{a\mid\sigma(n)\\a>(\log_2n)^2}}\frac1a\leq1.
\]
\end{lemma}

\begin{proof}
These are the consequences of Lemmas~2.1, 2.2 and~2.5 of
\cite{LucaPomerance} used below.  The first assertion is the valuation
conclusion established in the proof of Lemma~2.1(4).
\end{proof}

\begin{lemma}[Smooth-part counting input]\label{lem:smooth-part-input}
Let $Y\to\infty$.  For every fixed $c>0$, uniformly for $Y$-smooth
$d\leq Y^c$,
\[
 \#\{n\leq X:d\text{ is the largest $Y$-smooth divisor of }n\}
 \asymp \frac{X}{d\log Y}.
\]
Moreover, given $\eta>0$, $c$ may be chosen sufficiently large that the
integers $n\leq X$ whose largest $Y$-smooth divisor exceeds $Y^c$ account
for at most $\eta X$ integers for all sufficiently large $X$.
\end{lemma}

\begin{proof}
This is the smooth-part estimate and its accompanying tail bound used in
\cite[(3.1)]{LucaPomerance}.
\end{proof}

The following lemma selects a large subfamily with the regularity properties
needed in the collision estimate.

\begin{lemma}\label{lem:sv-regular}
There is a subfamily $\mathcal A(X)\subseteq\mathcal A_0(X)$ with
$\#\mathcal A(X)\gg_\delta X$ having the following properties.  For
$M=pqrk\in\mathcal A(X)$, let $d$ be the largest $Y$-smooth divisor of
$M$.  Then, for every $t\in\{k,\ell,m,M\}$,
\begin{enumerate}
\item $d=\gcd(t,\sigma(t))$, and $d$ is the largest $Y$-smooth divisor of
      both $t$ and $s(t)$;
\item $\sigma(t)/d$ is divisible by every prime at most $Y$, while every
      prime factor of $s(t)/d$ exceeds $y(t)$.
\end{enumerate}
In addition,
\[
 y(m)<\pi\leq m^{10/27},\quad \pi\text{ prime}
 \quad\Longrightarrow\quad \pi^2\nmid s(m),
\]
and
\begin{equation}\label{eq:sv-LP25}
  \sum_{\substack{a\mid\sigma(m)\\a>(\log_2X)^2}}\frac1a\leq1.
\end{equation}
There is also a constant $B_\delta>0$ such that
\begin{equation}\label{eq:sv-image-abundancy}
  \frac{\sigma(s(m))}{s(m)}\leq B_\delta.
\end{equation}
\end{lemma}

\begin{proof}
First discard the integers for which $k$ has a prime factor in
$(Y,y(X)]$.  Uniformly for $X^{1/120}<k\leq X^{1/60}$,
\[
 y(k)=y(X)+O\!\left(\frac1{\log_3X}\right),
\]
so
\[
 \sum_{Y<p\leq y(X)}\frac1p=o(1).
\]
Indeed,
\[
 \sum_{\substack{k\leq X^{1/60}\\
       p\mid k\text{ for some }Y<p\leq y(X)}}\frac1k
 \ll \log X\sum_{Y<p\leq y(X)}\frac1p=o(\log X).
\]
The discarded $k$ therefore account for $o(X)$ members of
$\mathcal A_0(X)$.

Apply Lemma~\ref{lem:LP-inputs} to each of $k,\ell,m,M$, with its
large-prime square-divisibility conclusion applied to $m$.

Each excluded set has asymptotic density zero.  These exclusions remove only
$o(X)$ members of $\mathcal A_0(X)$.  To justify this assertion, let
$\mathcal E\subseteq\N$ satisfy
$\#(\mathcal E\cap[1,T])=o(T)$.  Partial summation gives
\begin{equation}\label{eq:sv-harmonic-density-zero}
  \sum_{\substack{n\leq T\\n\in\mathcal E}}\frac1n=o(\log T).
\end{equation}
The separated factor ranges make the decompositions
$\ell=rk$, $m=q\ell$ and $M=pm$ unique.  An upper-bound estimate for the
prime $p$ then shows that the members for which one of $k,\ell,m$ lies in
$\mathcal E$ are
\[
  \ll\frac{X}{\log X}
       \sum_{\substack{t\leq X^{O(1)}\\t\in\mathcal E}}\frac1t=o(X),
\]
up to bounded reciprocal prime sums; the deletion at the level $M$ is
directly $o(X)$.

For every retained $t$, Lemma~\ref{lem:LP-inputs} and its first four
properties identify $\gcd(t,\sigma(t))$ with the largest $y(t)$-smooth
divisor of $t$ and give the corresponding divisibility properties of
$\sigma(t)$ and $s(t)$.  Since $Y\leq y(t)\leq y(X)$, the exclusion made in
the first paragraph and the fact that $p,q,r>y(X)$ show that the same integer
$d$ occurs at all four levels and prove (1)--(2).  The hypothesis
$q>m^{7/9}$ gives the displayed square-divisibility exclusion for $s(m)$,
and the final assertion of Lemma~\ref{lem:LP-inputs} gives
\eqref{eq:sv-LP25}.

After one further exclusion at the level $m$, justified by
\eqref{eq:sv-harmonic-density-zero}, \cite[Theorem~1.4]{Pollack} gives,
outside a set of asymptotic density zero,
\[
  \frac{s(s(m))}{s(m)}
  \leq \frac{s(m)}m+O\!\left((\log_2X)^{-1/4}\right).
\]
Together with \eqref{eq:sv-two-sided}, this yields
\eqref{eq:sv-image-abundancy}, for example with $B_\delta=C_\delta+2$ for all
sufficiently large $X$.
\end{proof}

Partition $\mathcal A(X)$ according to the integer $d$ in
Lemma~\ref{lem:sv-regular}.  The image sets belonging to distinct classes
are disjoint, because $d$ is also the largest $Y$-smooth divisor of $s(M)$.

\begin{lemma}\label{lem:sv-classes}
There are positive constants $c,c_1,c_2$, depending at most on $\delta$, and a set $\mathcal D$ of $Y$-smooth
integers $d\leq Y^c$ such that, writing
\[
  \mathcal A_d(X)=\{M\in\mathcal A(X):d
  \text{ is the largest $Y$-smooth divisor of }M\},
\]
one has, uniformly for $d\in\mathcal D$,
\begin{equation}\label{eq:sv-class-size}
  c_1\frac{X}{d\log Y}
  \leq\#\mathcal A_d(X)
  \leq c_2\frac{X}{d\log Y},
  \qquad
  \sum_{d\in\mathcal D}\frac1d\gg_\delta\log Y.
\end{equation}
The sets $s(\mathcal A_d(X))$, $d\in\mathcal D$, are pairwise disjoint.
\end{lemma}

\begin{proof}
Fix $\eta=\eta(\delta)>0$ with $\#\mathcal A(X)\geq\eta X$.
Apply Lemma~\ref{lem:smooth-part-input}, choosing $c$ so large that the
classes with $d>Y^c$ contain at most $\eta X/3$ integers.  Since
\[
  \sum_{P^+(d)\leq Y}\frac1d
  =\prod_{p\leq Y}\left(1-\frac1p\right)^{-1}
  \asymp\log Y,
\]
we may choose a sufficiently small $a=a(\delta)>0$ and let $\mathcal D$
consist of those $d\leq Y^c$ for which
\[
  \#\mathcal A_d(X)\geq a\frac{X}{d\log Y}.
\]
The classes outside $\mathcal D$ then contain at most $\eta X/3$ integers.
Thus the classes in $\mathcal D$ contain $\gg_\delta X$ integers.  The
uniform upper estimate above now gives
\[
  \frac{X}{\log Y}\sum_{d\in\mathcal D}\frac1d
  \gg_\delta X,
\]
which proves \eqref{eq:sv-class-size}.  The disjointness assertion follows
from Lemma~\ref{lem:sv-regular}.
\end{proof}

\subsection{The collision estimate}

\begin{proposition}\label{prop:sv-second-moment}
For $d\in\mathcal D$, put
\[
  R_d(u)=\#\{M\in\mathcal A_d(X):s(M)=u\}.
\]
Then, uniformly for $d\in\mathcal D$,
\begin{equation}\label{eq:sv-second-moment}
  \sum_uR_d(u)^2\ll_\delta\frac{X}{d\log Y}.
\end{equation}
\end{proposition}

\begin{proof}
We prove the collision estimate directly, following the organization of
\cite[(3.4), \S\S3.1--3.2]{LucaPomerance}.  All estimates needed for the
present family are recorded below.  We first note two reciprocal estimates.  Since $d$ is the largest $Y$-smooth divisor of every
member of $\mathcal A_d(X)$, the corresponding factor $k$ is $d$ times an
integer with no prime factor at most $Y$.  The elementary upper-bound sieve,
partial summation, and Mertens' theorem give
\begin{equation}\label{eq:sv-k-reciprocal}
 \sum_{k}\frac1k\ll\frac{\log X}{d\log Y}.
\end{equation}
Here and below a sum over a factor of a member of $\mathcal A_d(X)$ is
restricted to the ranges in \eqref{eq:sv-factorization}.  The reciprocal
prime sums over the fixed power intervals for $q$ and $r$ are bounded, so
\begin{equation}\label{eq:sv-m-reciprocal}
 \sum_m\frac1m\ll\frac{\log X}{d\log Y}.
\end{equation}

The diagonal contribution to $\sum_uR_d(u)^2$ is
$\#\mathcal A_d(X)$ and is bounded by \eqref{eq:sv-class-size}.  Consider
an off-diagonal collision, written as $M=pm$, $M'=p'm'$, with $m\ne m'$.
Then
\begin{equation}\label{eq:sv-collision}
 p\,s(m)+\sigma(m)=p'\,s(m')+\sigma(m').
\end{equation}
By Lemma~\ref{lem:sv-regular},
$\gcd(m,\sigma(m))=\gcd(m',\sigma(m'))=d$.  Put
\[
 \gcd(s(m),s(m'))=dh.
\]
Every prime factor of $h$ exceeds both $y(m)$ and $y(m')$.  Equation
\eqref{eq:sv-collision} shows that $dh$ divides
$\sigma(m)-\sigma(m')$.  This difference is nonzero.  Indeed, equality
would give $p s(m)=p's(m')$; for all sufficiently large $X$,
\[
 \min(p,p')>\tfrac12X^{8/15}>C_\delta X^{7/15}
 \geq\max\{s(m),s(m')\},
\]
which forces $p=p'$, then $s(m)=s(m')$, and finally $m=m'$.  Thus
\[
 A_1=\frac{s(m)}{dh},\qquad
 A_2=\frac{s(m')}{dh},\qquad
 A_3=\frac{|\sigma(m)-\sigma(m')|}{dh}
\]
are positive integers, with $\gcd(A_1,A_2)=1$.

For fixed $m,m'$, elementary linear Diophantine theory shows that there
are fixed integers $p_0,p'_0$ for which every integral solution has the form
\[
 p=p_0+A_2t,\qquad p'=p'_0+A_1t \qquad(t\in\mathbb Z).
\]
All relevant values of $t$ lie in an interval of length at most
\[
 T:=\frac{Xdh}{m s(m')}.
\]
The bounds
\[
 \lambda m'\leq s(m')\leq C_\delta m',
 \qquad mm'\leq X^{14/15},
\]
give $T\gg_\delta X^{1/15}$ and
$T\ll_\delta Xdh/(mm')$, so $\log T\asymp\log X$.  The affine forms are
primitive whenever a prime pair exists.  For example, a common divisor of
$p_0$ and $A_2$ would divide every value of $p=p_0+A_2t$, whereas
$p>A_2$ for all sufficiently large $X$; the second form is identical.
For odd primes the two primitive forms exclude at most two residue classes,
and the prime $2$ either leaves an allowed parity class or there are no
solutions.  The standard two-dimensional upper-bound sieve, in the form used in
\cite[(3.7)]{LucaPomerance}, therefore gives
\begin{equation}\label{eq:sv-two-prime-sieve}
 \ll_\delta
 \frac{Xdh}{mm'(\log X)^2}
 \frac{A_1}{\varphi(A_1)}
 \frac{A_2}{\varphi(A_2)}
 \frac{A_3}{\varphi(A_3)}
\end{equation}
possible prime pairs.

If $v\mid s(t)$, then
\begin{equation}\label{eq:sv-totient}
 \frac v{\varphi(v)}
 \leq\zeta(2)\frac{\sigma(v)}v
 \leq\zeta(2)\frac{\sigma(s(t))}{s(t)}
 \leq\zeta(2)B_\delta
 \qquad(t\in\{m,m'\}).
\end{equation}
This bounds the factors belonging to $A_1,A_2$, and also gives
$\varphi(h)\gg_\delta h$.

We now reduce the factor belonging to $A_3$.  Let $A_{3,1}$ contain its
prime factors at most $(\log_2X)^2$, let $A_{3,2}$ contain those in
$((\log_2X)^2,\log X]$, and let $A_{3,3}$ contain those exceeding
$\log X$.  Since $A_3\ll_\delta X^{7/15}$,
\[
 \frac{A_{3,1}}{\varphi(A_{3,1})}\ll\log Y,
 \qquad
 \frac{A_{3,3}}{\varphi(A_{3,3})}\ll1.
\]
Let $A'_{3,2}$ be the largest divisor of $A_{3,2}$ coprime to
$\sigma(m)$.  The primes removed from $A_{3,2}$ divide $\sigma(m)$ and
exceed $(\log_2X)^2$, so \eqref{eq:sv-LP25} gives
\[
 \frac{A_{3,2}}{\varphi(A_{3,2})}
 \ll\frac{A'_{3,2}}{\varphi(A'_{3,2})}.
\]
Consequently, it remains to prove
\begin{equation}\label{eq:sv-reduced-collision-sum}
 \frac{X\log Y}{(\log X)^2}
 \sum_{m\ne m'}
 \frac{dh}{mm'}
 \frac{A'_{3,2}}{\varphi(A'_{3,2})}
 \ll_\delta\frac{X}{d\log Y},
\end{equation}
where $dh=\gcd(s(m),s(m'))$.  Notice also that $A_3$ is divisible by every
prime at most $Y$.  Indeed, both $\sigma(m)/d$ and $\sigma(m')/d$ have this
property by Lemma~\ref{lem:sv-regular}, while every prime factor of $h$
exceeds $Y$.

Suppose first that $h>X^{10/33}$.  Write $m=q\ell$ and $m'=q'\ell'$.  We claim that
$\gcd(h,s(\ell)\sigma(\ell))=1$.  If a prime $\pi$ divided both $h$
and $\sigma(\ell)$, then the identity
$s(m)=q s(\ell)+\sigma(\ell)$ would imply either $\pi=q$ or
$\pi\mid s(\ell)$.  In the first case,
$\pi\mid\gcd(m,\sigma(m))=d$; in the second case,
$\pi\mid\ell$ and $\pi\mid\sigma(\ell)$, so again $\pi\mid d$.
Both conclusions contradict the fact that every prime factor of $h$ exceeds
$Y$, whereas $d$ is $Y$-smooth.  If instead $\pi\mid h$ and
$\pi\mid s(\ell)$, then the same identity gives
$\pi\mid\sigma(\ell)$, which has just been excluded.  Thus both
$s(\ell)$ and $\sigma(\ell)$ are units modulo $h$.  It follows that
$q$ is fixed in a reduced residue class modulo $h$, and the same holds for
$q'$.  Moreover,
\[
 m\equiv\sigma(m)\equiv\sigma(m')\equiv m'\pmod h.
\]
Eliminating $q,q'$ gives
\begin{equation}\label{eq:sv-large-h-congruence}
 s(\ell')\ell\sigma(\ell)-s(\ell)\ell'\sigma(\ell')
 \equiv0\pmod h.
\end{equation}
The absolute value of the left side is at most
\[
 2C_\delta(C_\delta+1)\max(\ell,\ell')^3
 \ll_\delta X^{3/10}=o(X^{10/33}),
\]
so it vanishes.  For every prime power $\pi^a\Vert d$,
Lemma~\ref{lem:sv-regular} gives
$\pi^a\Vert s(\ell)$ and $\pi^{a+1}\mid\sigma(\ell)$; it also excludes
primes exceeding $Y$ from $\gcd(\ell,s(\ell))$.  Hence
\[
 \gcd(\ell^2,s(\ell))=d,
 \qquad
 \gcd((\ell')^2,s(\ell'))=d.
\]
The equality in \eqref{eq:sv-large-h-congruence} can be written as
\[
 \frac{\ell^2}{s(\ell)}+\ell
 =\frac{(\ell')^2}{s(\ell')}+\ell'.
\]
After writing $\ell=dL$ and $\ell'=dL'$, both fractions are reduced with
denominators $s(\ell)/d$ and $s(\ell')/d$.  Their difference is an integer,
so the denominators are equal.  Thus $s(\ell)=s(\ell')$, and the displayed
equality, viewed as a strictly increasing quadratic in $\ell$, gives
$\ell=\ell'$.

In this range Mertens' theorem gives
$A'_{3,2}/\varphi(A'_{3,2})\ll\log_2X/\log Y$.  Summing the left side of
\eqref{eq:sv-reduced-collision-sum} and using $\ell=\ell'$ gives
\[
 \ll\frac{Xd\log_2X}{(\log X)^2}
 \sum_{q,q',\ell,h}\frac{h}{qq'\ell^2}.
\]
For $q>q'$, the congruence classes just obtained give
$\sum_q1/q\ll(\log X)/h$, even if primality is discarded.  Summing next
over $h\mid s(q'\ell)$ is bounded by
$\tau(s(q'\ell))=X^{o(1)}$.  The reciprocal sum over $q'$ is bounded and
$\sum_\ell\ell^{-2}\ll X^{-1/15}$.  The contribution is therefore
\[
 dX^{14/15+o(1)}=o\!\left(\frac{X}{d\log Y}\right),
\]
uniformly for $d\leq Y^c$.

It remains to treat $h\leq X^{10/33}$.  Let $A''_{3,2}$ be the largest
divisor of $A'_{3,2}$ coprime to $s(m')$.  Equation
\eqref{eq:sv-totient} gives
\[
 \frac{A'_{3,2}}{\varphi(A'_{3,2})}
 \ll_\delta\frac{A''_{3,2}}{\varphi(A''_{3,2})}.
\]
We first note that $h$ is squarefree in this range.  Indeed, if
$\pi^2\mid h$, then $\pi^2\mid s(m)$ and $\pi>y(m)$.  By
Lemma~\ref{lem:LP-inputs}, this forces $\pi>m^{10/27}$.  Since
$m>X^{17/40}$,
\[
 \pi^2>m^{20/27}>X^{17/54}>X^{10/33}\geq h,
\]
a contradiction.

Fix $m',h,k$.  The congruences
\[
 qrk\equiv(q+1)(r+1)\sigma(k)\equiv m'\pmod h
\]
may therefore be solved for the required symmetric functions.  Indeed, any prime
dividing both $h$ and either $m$ or $\sigma(m)$ would divide
$\gcd(m,\sigma(m))=d$, which is impossible because every prime factor of
$h$ exceeds $Y$.  Thus $k$ and $\sigma(k)$ are units modulo $h$.  The two
congruences therefore
determine $qr$ and $q+r$ modulo $h$.  Since $h$ is squarefree, there are at
most $\tau(h)$ possible ordered residue pairs
\[
 q\equiv a\pmod h,\qquad r\equiv b\pmod h.
\]
For one such pair put
\[
 \mathfrak f(qr)=
 \sum_{\substack{(\log_2X)^2<\pi\leq\log X\\
        \pi\mid\sigma(kqr)-\sigma(m')\\
        \pi\nmid h\sigma(kqr)}}\frac1\pi,
\]
where $\pi$ is prime.  Mertens' theorem gives
\begin{equation}\label{eq:sv-intermediate-totient}
 \frac{A''_{3,2}}{\varphi(A''_{3,2})}
 \ll_\delta 1+\frac{\log_2X}{\log Y}\mathfrak f(qr).
\end{equation}
Indeed, when $\mathfrak f(qr)\leq1$, the logarithm of the left side is
$O(\mathfrak f(qr))$.  When $\mathfrak f(qr)>1$, Mertens' theorem bounds
the contribution of the full prime range by
$O(\log_2X/\log Y)$.

For fixed $m',h,k,r$ and a prime $\pi$ occurring in
$\mathfrak f(qr)$, the congruence
\[
 (q+1)\sigma(kr)\equiv\sigma(m')\pmod\pi
\]
fixes $q$ modulo $\pi$.  The condition
$\pi\nmid\sigma(kqr)$ ensures that $\sigma(kr)$ is invertible modulo
$\pi$.  Since $\pi\nmid h$, this congruence combines with
$q\equiv a\pmod h$.  Brun--Titchmarsh and partial summation, followed by
\eqref{eq:sv-totient}, give
\begin{equation}\label{eq:sv-f-sum}
 \sum_{q,r}\frac{\mathfrak f(qr)}{qr}
 \ll_\delta
 \begin{cases}
  \dfrac{\log X}{hX^{1/20}},&X^{1/20}<h\leq X^{10/33},\\[6pt]
  \dfrac{1}{h^2(\log_2X)^2},&h\leq X^{1/20}.
 \end{cases}
\end{equation}
Here the sums are over the fixed residue pair.  To see
\eqref{eq:sv-f-sum}, the reciprocal prime sum over $q$ in its progression
is $O(1/\varphi(\pi h))\ll_\delta1/(\pi h)$.  Summing the extra weight
$1/\pi$ gives $O(1/(h(\log_2X)^2))$.  The reciprocal sum over $r$ in its
class is
\[
 \ll_\delta
 \begin{cases}
  (\log X)X^{-1/20},&h>X^{1/20},\\
  h^{-1},&h\leq X^{1/20},
 \end{cases}
\]
which proves the estimate.  The same progression bounds without
$\mathfrak f$ and \eqref{eq:sv-intermediate-totient} yield
\begin{equation}\label{eq:sv-qr-sum}
 \sum_{q,r}\left(\frac1{qr}
 +\frac{\log_2X}{\log Y}\frac{\mathfrak f(qr)}{qr}\right)
 \ll_\delta
 \begin{cases}
  \dfrac{\log X\log_2X}{hX^{1/20}\log Y},
       &X^{1/20}<h\leq X^{10/33},\\[6pt]
  h^{-2},&h\leq X^{1/20}.
 \end{cases}
\end{equation}

Using \eqref{eq:sv-k-reciprocal}, the contribution for fixed $m',h$ to
\eqref{eq:sv-reduced-collision-sum} is at most
\[
 \frac{X\log Y}{(\log X)^2}\frac{dh}{m'}
 \frac{\log X}{d\log Y}\tau(h)
 \times\text{the applicable bound in \eqref{eq:sv-qr-sum}}.
\]
For $X^{1/20}<h\leq X^{10/33}$ this is
\[
 \ll_\delta
 \frac{X^{19/20}\log_2X}{\log Y}\frac{\tau(h)}{m'}.
\]
Now $\sum_{h\mid s(m')}\tau(h)\leq\tau(s(m'))^2=X^{o(1)}$, and
\eqref{eq:sv-m-reciprocal} shows that the total is
\[
 \frac{X^{19/20+o(1)}}d
 =o\!\left(\frac{X}{d\log Y}\right).
\]
For $h\leq X^{1/20}$ the fixed-$m',h$ contribution is
\[
 \ll_\delta\frac{X}{\log X}\frac{\tau(h)}{hm'}.
\]
Finally,
\begin{equation}\label{eq:sv-divisor-sum}
 \sum_{h\mid s(m')}\frac{\tau(h)}h
 \leq\left(\frac{\sigma(s(m'))}{s(m')}\right)^2
 \leq B_\delta^2.
\end{equation}
Summing over $h$ and then using \eqref{eq:sv-m-reciprocal} gives
$O_\delta(X/(d\log Y))$.  Together with the large-$h$ and diagonal
contributions, this proves \eqref{eq:sv-second-moment}.
\end{proof}

\begin{proof}[Proof of Theorem~\ref{thm:small-values}]
For each $d\in\mathcal D$, Cauchy--Schwarz and
Lemma~\ref{lem:sv-classes} and Proposition~\ref{prop:sv-second-moment} give
\[
  \#s(\mathcal A_d(X))
  \geq\frac{\#\mathcal A_d(X)^2}{\sum_uR_d(u)^2}
  \gg_\delta\frac{X}{d\log Y}.
\]
The image sets are disjoint, and hence
\[
  \#s(\mathcal A(X))
  \gg_\delta\frac{X}{\log Y}
       \sum_{d\in\mathcal D}\frac1d
  \gg_\delta X.
\]
If $N=s(M)$ with $M\in\mathcal A(X)$, then
\eqref{eq:sv-basic} and \eqref{eq:sv-two-sided} give
\[
  f(N)\leq M<\delta N.
\]
Moreover, $N=s(M)<C_\delta M\leq C_\delta X$.  Replacing $X$ by $X/C_\delta$ proves the theorem
for $0<\delta\leq1$, and the case $\delta>1$ follows from $\delta=1$.
\end{proof}

\subsection{Consequences for empirical distributions}

For $x\geq0$, let $\delta_x$ denote the unit point mass at $x$.  For every
$X$ with $\Rcal\cap[1,X]\ne\varnothing$, define the probability measure
\begin{equation}\label{eq:empirical-measures}
 \nu_X=\frac1{\#(\Rcal\cap[1,X])}
 \sum_{\substack{N\leq X\\N\in\Rcal}}\delta_{f(N)/N}.
\end{equation}
A family of probability measures $\{\nu_X\}$ on $[0,\infty)$ is called
\emph{tight} if
\[
 \lim_{T\to\infty}\sup_X\nu_X((T,\infty))=0.
\]

\begin{corollary}\label{cor:limit-law-tails}
Let $X_j\to\infty$, and suppose that $\nu_{X_j}$ converges weakly to a
probability measure $\nu$ on $[0,\infty)$.  Then
\[
 0\in\supp(\nu),\qquad \nu(\{0\})=0.
\]
Moreover, for every $\eps>0$ there is a constant $c_\eps>0$ such that
\[
 \nu([A,\infty))\geq c_\eps A^{-1-\eps}
 \qquad(A\geq1).
\]
In particular, $\nu$ has unbounded support.
\end{corollary}

\begin{proof}
By \eqref{eq:exact-representability},
$\#(\Rcal\cap[1,X])=\lfloor X\rfloor-2$ for $X\geq5$.
Theorem~\ref{thm:almost-log-tail} gives, as $A\to\infty$,
\[
 \liminf_{X\to\infty}\nu_X([A,\infty))
 \geq\left(\frac{e^{-\gamma}}2-o(1)\right)
       \frac{\log_3A}{(\log A)(\log_2A)}.
\]
For every fixed $\eps>0$ this is $\gg_\eps A^{-1-\eps}$ for all
sufficiently large $A$; monotonicity in $A$ absorbs the remaining bounded
range into the implied constant.  Since $[A,\infty)$ is closed, the
portmanteau theorem yields the displayed bound in the corollary for
$\nu([A,\infty))$.  Thus the support is unbounded.

For every fixed $\delta>0$, Theorem~\ref{thm:small-values} gives
\[
 \liminf_{X\to\infty}\nu_X([0,\delta])>0.
\]
Again $[0,\delta]$ is closed, so $\nu([0,\delta])>0$.  Every neighborhood
of $0$ therefore has positive $\nu$-measure, and $0\in\supp(\nu)$.

Finally, Theorem~\ref{thm:small-upper} gives constants $C,c>0$ such that,
for all sufficiently small $\delta>0$,
\[
 \limsup_{j\to\infty}\nu_{X_j}([0,\delta])
 \leq C\exp\!\left\{-\exp((1/\delta)^c)\right\}.
\]
Hence
\[
 \limsup_{j\to\infty}\nu_{X_j}([\delta,\infty))
 \geq1-C\exp\!\left\{-\exp((1/\delta)^c)\right\}.
\]
The set $[\delta,\infty)$ is closed, so the portmanteau theorem gives
\[
 \nu([0,\delta))
 \leq C\exp\!\left\{-\exp((1/\delta)^c)\right\}.
\]
Letting $\delta\downarrow0$ proves $\nu(\{0\})=0$.
\end{proof}

\begin{remark}
The preceding corollary is conditional on the existence of a weakly
convergent subsequence.  The large- and small-value theorems alone do not
establish tightness of the family $\{\nu_X\}$ or convergence to a limiting
distribution.
\end{remark}

\section{Further structure and the density-one question}\label{sec:tightness}

Recall the divisor-prefix sums $\sigma_j(m)$ defined in
Section~\ref{sec:introduction}.
For $A\geq1$ and $N\in\N$, define
\[
 r_A(N)=\#\{(m,j):\sigma_j(m)=N,\ m\leq AN\}.
\]
Equivalently, $r_A(N)$ counts the pairs $m=ed\leq AN$ for which
$N=F_e(d)$, with endpoint divisor $d=d_{\tau(m)-j}(m)$.
Consider the density-one statement
\begingroup
\renewcommand{\theequation}{T}
\begin{equation}\label{eq:T}
 \lim_{A\to\infty}\limsup_{X\to\infty}
 \frac1X\#\{N\leq X:r_A(N)=0\}=0.
\end{equation}
\endgroup
Here $A$ is fixed while $X\to\infty$, and only then is $A$ allowed to
tend to infinity.  The first result below shows that the first moment of
$r_A$ tends to infinity within every fixed residue class.  The remaining
obstruction is collision: many representation pairs may correspond to the
same represented value.

\subsection{A first moment in every residue class}

For fixed $Q\geq1$ and $\rho\pmod Q$, put
\[
 \underline\lambda_{Q,\rho}(A)
 =\liminf_{X\to\infty}\frac QX
 \sum_{\substack{N\leq X\\N\equiv\rho\pmod Q}}r_A(N).
\]

\begin{proposition}[First moment in a fixed residue class]
\label{prop:class-first-moment}
For every fixed $Q\geq1$ and every residue class $\rho\pmod Q$,
\[
 \underline\lambda_{Q,\rho}(A)
 \gg_{Q,\rho}\frac{A}{(\log A)^2}
 \qquad(A\to\infty).
\]
In particular, this lower class mean tends to infinity in every residue
class.
\end{proposition}

\begin{proof}
Fix $Q$ and $\rho$, and set $d_0=\gcd(\rho,Q)$.  We first choose a fixed
integer $h$ satisfying
\begin{equation}\label{eq:h-gcd}
 \gcd(\sigma(h),Q)=d_0.
\end{equation}
For each prime power $\ell^k\Vert d_0$, Dirichlet's theorem on primes in
arithmetic progressions permits the choice of a different auxiliary prime
$q_\ell$ congruent to $1$ modulo every prime dividing $Q$.  When $\ell=2$,
we also require $q_2\equiv1\pmod4$.  Put
\[
 h_\ell=q_\ell^{\ell^k-1},\qquad
 h=\prod_{\ell^k\Vert d_0}h_\ell,
\]
with $h=1$ when $d_0=1$.  The lifting-the-exponent formula gives
$v_\ell(\sigma(h_\ell))=k$.  If $r$ is a prime divisor of $Q$ with $r\ne\ell$, then
$\sigma(h_\ell)\equiv\ell^k\not\equiv0\pmod r$.  Multiplicativity of
$\sigma$ proves \eqref{eq:h-gcd}.

The elementary criterion for a linear congruence now gives a unit
$u_\rho\pmod Q$ for which
\begin{equation}\label{eq:u-class}
 \sigma(h)u_\rho\equiv\rho\pmod Q.
\end{equation}
This includes the class $\rho=0$.  Prime-power by prime-power, divide by the
common valuation in \eqref{eq:h-gcd}, solve for a unit modulo the remaining
power, and combine the solutions by the Chinese remainder theorem.

Dirichlet's theorem also supplies infinitely many primes $e$ satisfying
\[
 e>\max\{2,h,P^+(Q)\},\qquad e\equiv-1\pmod Q.
\]
For each such prime $e$, define
\[
 \mathcal S_e(u_\rho)=\{t\geq1:t\equiv u_\rho\pmod Q,\
 (t,\mathfrak Q(e))=1\}.
\]
The Chinese remainder theorem shows that this periodic set has density
\begin{equation}\label{eq:delta-e}
 \delta_e=\frac1Q
 \prod_{\substack{p\leq e\\p\nmid Q}}\left(1-\frac1p\right).
\end{equation}
Every prime factor of $t\in\mathcal S_e(u_\rho)$ is larger than $e$.

For $t\in\mathcal S_e(u_\rho)$, the divisors of $eht$ below $e$ are exactly
the divisors of $h$.  Complementary divisors therefore give
\begin{equation}\label{eq:He}
 F_e(ht)=\sigma(h)H_e(t),
 \qquad H_e(t)=(e+1)\sigma(t)-et=t+(e+1)s(t),
\end{equation}
where $s(t)=\sigma(t)-t$.  Since $Q\mid e+1$, equations
\eqref{eq:u-class} and \eqref{eq:He} give
\[
 F_e(ht)\equiv\sigma(h)t\equiv\rho\pmod Q.
\]
The corresponding representing integer is $m=eht$.  Since $H_e(t)\geq t$,
\begin{equation}\label{eq:representing-ratio}
 \frac{m}{F_e(ht)}\leq\frac{eh}{\sigma(h)}\leq eh.
\end{equation}

We next bound the represented value.  Complementary divisors give
\[
 \frac{H_e(t)}t
 =1+(e+1)\sum_{\substack{a\mid t\\a>1}}\frac1a.
\]
For a fixed integer $a>1$ with $(a,\mathfrak Q(e))=1$, the Chinese remainder theorem
shows that the set of $t\in\mathcal S_e(u_\rho)$ divisible by $a$ has density
$\delta_e/a$.  Since $\sum_{a\geq1}a^{-2}$ converges, dominated
convergence permits the interchange of the limit and the divisor sum.
Therefore
\begin{align*}
 \lim_{T\to\infty}\frac1T
 \sum_{\substack{t\leq T\\t\in\mathcal S_e(u_\rho)}}\frac{H_e(t)}t
 &=\delta_e\left(1+(e+1)
 \sum_{\substack{a>1\\(a,\mathfrak Q(e))=1}}\frac1{a^2}\right)\\
 &\leq3\delta_e.
\end{align*}
Every integer $a>1$ coprime to $\mathfrak Q(e)$ exceeds $e$.  Hence
\[
 (e+1)\sum_{\substack{a>1\\(a,\mathfrak Q(e))=1}}\frac1{a^2}
 \leq(e+1)\sum_{a>e}\frac1{a^2}
 \leq\frac{e+1}{e}\leq\frac32,
\]
which proves the last inequality.  Markov's
inequality shows that
\[
 \mathcal G_e=\{t\in\mathcal S_e(u_\rho):H_e(t)\leq6t\}
\]
has lower density at least $\delta_e/2$.

For a fixed $A$, restrict to primes $e$ satisfying
\[
 \max\{2,h,P^+(Q)\}<e\leq A/h,
 \qquad e\equiv-1\pmod Q.
\]
If $t\in\mathcal G_e$ and $t\leq X/(6\sigma(h))$, then the represented value in
\eqref{eq:He} is at most $X$, lies in the class $\rho$, and has a representing integer
$m\leq AN$ by \eqref{eq:representing-ratio}.  Distinct pairs $(e,t)$ give
distinct representing integers.  Indeed, if $e<e'$ and $eht=e'ht'$, then $e\mid t'$,
contrary to the fact that every prime factor of $t'$ exceeds $e'$.  If
$e=e'$, equality of the representing integers forces $t=t'$.  Hence these pairs give distinct representation pairs $(m,j)$, even if some
represented values coincide.  For each eligible prime $e$, the number of
such $t\leq X/(6\sigma(h))$ is at least
$(\delta_e/2+o(1))X/(6\sigma(h))$.  Multiplying by $Q/X$ and summing over
$e$ gives
\begin{equation}\label{eq:class-lower}
 \underline\lambda_{Q,\rho}(A)
 \geq\frac1{12\sigma(h)}
 \sum_{\substack{\max\{2,h,P^+(Q)\}<e\leq A/h\\
                  e\text{ prime}\\e\equiv-1\pmod Q}}
 \prod_{\substack{p\leq e\\p\nmid Q}}\left(1-\frac1p\right).
\end{equation}
Mertens' product theorem gives
\[
 \prod_{\substack{p\leq e\\p\nmid Q}}
 \left(1-\frac1p\right)
 \sim\frac Q{\varphi(Q)}\frac{e^{-\gamma}}{\log e}.
\]
For $Z\geq2$, the prime number theorem in arithmetic progressions, followed
by partial summation, gives, as $Z\to\infty$,
\[
 \sum_{\substack{e\leq Z\\e\text{ prime}\\e\equiv-1\pmod Q}}
 \frac1{\log e}
 \sim\frac{Z}{\varphi(Q)(\log Z)^2}.
\]
The omitted primes below the fixed threshold $\max\{2,h,P^+(Q)\}$ do not
change this asymptotic.  Taking $Z=A/h$ in \eqref{eq:class-lower} proves the
proposition.
\end{proof}

\begin{remark}
Proposition~\ref{prop:class-first-moment} counts representation pairs.  A small set
of represented values could carry many pairs, so the proposition alone gives no bound
for $\#\{N\leq X:r_A(N)=0\}$.  Thus the proposition does not by itself imply
\eqref{eq:T}.
\end{remark}

\subsection{The range of the sum-of-divisors function}

\begin{lemma}[The range of the full divisor sum has density zero]
\label{lem:sigma-range-zero}
The set $\sigma(\N)=\{\sigma(m):m\geq1\}$ has asymptotic density zero.
\end{lemma}

\begin{proof}
If $N=\sigma(m)\leq X$, then $m\leq X$.  For every fixed integer $q\geq2$,
the values divisible by $q$ number at most $X/q$, while all remaining values
arise from $o_q(X)$ exceptional sources by
Lemma~\ref{lem:fixed-modulus-normality}.  Hence
\[
 \upperdens(\sigma(\N))\leq\frac1q.
\]
Letting $q\to\infty$ proves the claim.
\end{proof}

\subsection{The forced modulus}\label{subsec:forced-modulus}
Lemma~\ref{lem:bounded} shows that a value $N=F_e(d)$ agrees with
$\sigma(ed)$ modulo some prime below $e\Lambda_e$.  The modulus can be
written down exactly.  Fix $e\geq2$ and $d\geq1$, put $m=ed$, and let
\begin{equation}\label{eq:fm-signature}
 D=\{j:1\leq j<e,\ j\mid m\}
\end{equation}
be the set of divisors of $m$ below $e$.  Put
\[
 M=\lcm D,\qquad C=\sum_{j\in D}\frac Mj,\qquad g=\gcd(e,M),\qquad
 K_{e,d}=\frac eg\,C .
\]
Since $1\in D$, all four are positive integers.  For $1\leq j<e$ write
$h=\gcd(j,e)$.  Then
\[
 j\mid ed\quad\Longleftrightarrow\quad \frac jh\mid\frac eh\,d
 \quad\Longleftrightarrow\quad \frac jh\mid d,
\]
because $j/h$ and $e/h$ are coprime.  Thus $D$ depends only on the residue
of $d$ modulo $\lcm\{j/\gcd(j,e):1\leq j<e\}$, and for fixed $e$ the
integer $K_{e,d}$ takes only finitely many values.
\begin{lemma}[Exact forced modulus]\label{lem:fm-modulus}
With the notation above,
\begin{equation}\label{eq:fm-reflection}
 F_e(d)=\sigma(m)-\frac mM\,C ,
\end{equation}
the integer $M/g$ divides $d$, and
\begin{equation}\label{eq:fm-congruence}
 F_e(d)\equiv\sigma(m)\pmod{K_{e,d}} .
\end{equation}
Moreover $K_{e,d}$ is the largest modulus for which
\eqref{eq:fm-congruence} holds for every $d$ with the same set $D$.
\end{lemma}
\begin{proof}
The map $r\mapsto m/r$ is an involution of the divisors of $m$, and $r>d$
exactly when $m/r<e$.  Hence
\[
 \sigma(m)-F_e(d)=\sum_{\substack{r\mid m\\ r>d}}r
 =\sum_{j\in D}\frac mj=\frac mM\sum_{j\in D}\frac Mj=\frac mM\,C .
\]
Every $j\in D$ divides $m$, so $M\mid m$ and $m/M$ is an integer.  This
proves \eqref{eq:fm-reflection}.
Since $g=\gcd(e,M)$, the integers $M/g$ and $e/g$ are coprime.  From
$M\mid ed$ we get $(M/g)\mid(e/g)d$, hence $(M/g)\mid d$.  Therefore
\[
 \frac mM=\frac{ed}M=\frac eg\cdot\frac d{M/g}
\]
is $e/g$ times an integer, and by \eqref{eq:fm-reflection} the difference
$\sigma(m)-F_e(d)$ is a multiple of $(e/g)C=K_{e,d}$.
For the last assertion take $d_0=M/g$, so that $ed_0=\lcm(e,M)$.  Every
$j\in D$ divides $M$, hence $ed_0$.  Conversely $ed_0$ divides $ed$, so
every divisor of $ed_0$ below $e$ lies in $D$.  Thus $d_0$ has the same set
$D$, and at $d_0$ the difference $(ed_0/M)C$ equals $(e/g)C$ exactly.  Any
modulus valid for every $d$ with this set $D$ divides that difference,
hence divides $K_{e,d}$.
\end{proof}
If $P^-(d)\geq e$, then $D$ is the set of divisors of $e$ below $e$.
Indeed, for $j<e$ the integer $j/\gcd(j,e)$ is below $e\leq P^-(d)$, so it
divides $d$ only if it equals $1$, that is, only if $j\mid e$.  In this
case $M\mid e$, $g=M$, and
\[
 K_{e,d}=\sum_{\substack{j\mid e\\ j<e}}\frac ej=\sigma(e)-1 .
\]
So on $e$-rough inputs, $F_e(d)\equiv\sigma(ed)\pmod{\sigma(e)-1}$.
For prime cofactors the modulus always contains the cofactor itself, which
bounds the density of each prime range.
\begin{corollary}[Prime cofactors]\label{cor:fm-prime-ceiling}
Let $p$ be prime.  Then $F_p(d)\equiv\sigma(pd)\pmod p$ for every $d\geq1$,
and
\[
 \#\bigl(F_p(\N)\cap[1,X]\bigr)\leq\frac Xp+o_p(X),
 \qquad\text{so}\qquad
 \upperdens\bigl(F_p(\N)\bigr)\leq\frac1p .
\]
\end{corollary}
\begin{proof}
Every $j\in D$ satisfies $j<p$, so $p\nmid M$ and $g=1$.  Hence
$p\mid K_{p,d}$, and the congruence follows from \eqref{eq:fm-congruence}.
Now let $N=F_p(d)\leq X$ with $p\nmid N$.  Then $p\nmid\sigma(pd)$, and
$pd\leq pN\leq pX$ because $d\leq F_p(d)$.  Distinct $N$ have distinct
sources $pd$, so these $N$ number at most
$\#\{m\leq pX:p\nmid\sigma(m)\}$, which is $o_p(X)$ by
Lemma~\ref{lem:fixed-modulus-normality}.  Every other value of $F_p$ up to
$X$ is a multiple of $p$.
\end{proof}
For $p=2$ the bound $1/2$ is weaker than Proposition~\ref{prop:theta-two}
below.  What the corollary adds is uniformity in $p$.  Each prime range on
its own has density at most $1/p$, so a union of density close to one has
to draw on many cofactors at once.

\subsection{Bounded-cofactor coverage}

For a positive integer $A$, put
\[
 G_A=\bigcup_{e\leq A}F_e(\N),\qquad
 \Lambda_A=\lcm(1,\ldots,A),\qquad P_A=A\Lambda_A.
\]
Every $N=F_e(d)$ with
$e\leq A$ has a representing integer $m=ed\leq AN$, because $F_e(d)\geq d$.
Thus
\begin{equation}\label{eq:GA-sufficient}
 G_A\subseteq\{N:r_A(N)>0\}.
\end{equation}

\begin{corollary}[Every fixed cofactor union has a defect]
\label{cor:fixed-cofactor-defect}
For every fixed $A\geq1$,
\[
 \lowerdens(\N\setminus G_A)>0.
\]
More precisely, as $A\to\infty$,
\begin{equation}\label{eq:fixed-cofactor-defect}
 \lowerdens\{N:N\text{ squarefree},\ N\notin G_A\}
 \geq\left(\frac{e^{-\gamma}}2-o(1)\right)
       \frac{\log_3A}{(\log A)(\log_2A)}.
\end{equation}
\end{corollary}

\begin{proof}
If $N\in\Rcal$ and $f(N)>AN$, then $N\notin G_A$.  The quantitative
assertion follows from Theorem~\ref{thm:almost-log-tail} with $T=A$.
For any fixed smaller $A$, apply that theorem with one sufficiently large
fixed $T\geq A$.
\end{proof}

Write
\[
 \delta_A=\frac{\varphi(\mathfrak Q(P_A))}{\mathfrak Q(P_A)},\qquad
 \eta_A=\upperdens\{N:(N,\mathfrak Q(P_A))>1,\ N\notin G_A\}.
\]
For $A\geq2$, Lemmas~\ref{lem:sigma-rate} and~\ref{lem:bounded}, with
$E=A$ and $P=P_A$, give
\[
 \#\{N\leq X:(N,\mathfrak Q(P_A))=1,\ N\in G_A\}=o_A(X).
\]
Thus $G_A$ has density zero among the integers coprime to
$\mathfrak Q(P_A)$.  Hence the coprime part of $\N\setminus G_A$ has density
$\delta_A$.  Since the coprime and noncoprime parts are disjoint,
\begin{equation}\label{eq:bounded-cofactor-complement}
 1-\lowerdens(G_A)=\delta_A+\eta_A.
\end{equation}
Moreover, Mertens' theorem and the prime number theorem, in the form
$\psi(A)=A+o(A)$, give
$\log P_A=\log A+\psi(A)=A+o(A)$ and hence
\begin{equation}\label{eq:delta-A-asymptotic}
 \delta_A=\frac{e^{-\gamma}+o(1)}A.
\end{equation}

The two terms in \eqref{eq:bounded-cofactor-complement} are of different
orders, and it is $\eta_A$ that dominates.  Write
\[
 L(A)=\frac{\log_3A}{(\log A)(\log_2A)} .
\]

\begin{proposition}\label{prop:fm-eta}
As $A\to\infty$,
\[
 \eta_A\geq\Bigl(\frac{e^{-\gamma}}2-o(1)\Bigr)L(A),
 \qquad\text{so}\qquad A\eta_A\to\infty .
\]
\end{proposition}

\begin{proof}
By \eqref{eq:fixed-cofactor-defect},
$1-\lowerdens(G_A)\geq(e^{-\gamma}/2-o(1))L(A)$.  Since $AL(A)\to\infty$,
\eqref{eq:delta-A-asymptotic} gives $\delta_A=o(L(A))$.  Subtracting in
\eqref{eq:bounded-cofactor-complement} proves the first claim, and the
second follows from $AL(A)\to\infty$.
\end{proof}

One might guess that coprimality to $\mathfrak Q(P_A)$ is asymptotically
the only obstruction, so that $1-\lowerdens(G_A)\sim e^{-\gamma}/A$.  The
proposition rules this out.  The integers missed by $G_A$ for other reasons
outnumber the coprime ones by a factor that tends to infinity.

Now set 
\[
 \theta_2=\upperdens\bigl(F_2(\N)\bigr).
\]

\begin{proposition}\label{prop:theta-two}
One has
\[
 \theta_2\leq\frac12-0.0602757=0.4397243.
\]
\end{proposition}

\begin{proof}
The only divisor of $2d$ exceeding $d$ is $2d$, and hence
\begin{equation}\label{eq:F2-aliquot}
 F_2(d)=\sigma(2d)-2d=s(2d).
\end{equation}
Thus $F_2(\N)\subseteq s(\N)$ and the cofactor-two range omits every
untouchable number.  Write
\[
 \mathcal U=\N\setminus s(\N).
\]

We first note that the odd part of $F_2(\N)$ is negligible.  If $m$ is even,
then $s(m)$ is odd precisely when $\sigma(m)$ is odd, which occurs precisely
when $m$ is a square or twice a square.  For even $m\geq4$ one also has
$s(m)\geq1+m/2$; the exceptional case $m=2$ is harmless.  Thus $s(m)\leq X$
implies $m<2X$, apart from at most this one value.  Consequently,
\begin{equation}\label{eq:F2-odd-negligible}
 \#\{N\leq X:N\text{ odd},\ N\in F_2(\N)\}=O(\sqrt X).
\end{equation}

Almost all untouchable numbers are even.  Indeed, let $N$ be odd.  If
$N-1=p+q$ with distinct primes $p$ and $q$, then the proper divisors of $pq$
are $1,p,q$, and therefore $s(pq)=N$.  The exceptional-set theorem of
Montgomery and Vaughan shows that the even integers up to $X$ that are not a
sum of two primes number $O(X^{1-\kappa})$ for some $\kappa>0$; the diagonal
possibilities $N-1=2p$ contribute $O(X/\log X)$ more \cite{MVGoldbach}.
Hence
\begin{equation}\label{eq:odd-untouchables}
 \#\{N\leq X:N\text{ odd},\ N\in\mathcal U\}=o(X).
\end{equation}
Chen and Zhao proved \cite{ChenZhao} that
\[
 \lowerdens(\mathcal U)\geq0.0602757.
\]
By \eqref{eq:odd-untouchables}, the number of even
untouchables not exceeding $X$ is at least
$(0.0602757+o(1))X$.  There are $X/2+O(1)$ even integers up to $X$, and by
\eqref{eq:F2-odd-negligible} all but $O(\sqrt X)$ elements of
$F_2(\N)\cap[1,X]$ are even.  Since \eqref{eq:F2-aliquot} shows that this
range misses every member of $\mathcal U$, we obtain
\[
 \#\bigl(F_2(\N)\cap[1,X]\bigr)
 \leq\frac X2-0.0602757X+o(X).
\]
Taking the upper limit proves the proposition.
\end{proof}

\begin{corollary}\label{cor:eta-two}
For the residual density defined above,
\[
 \eta_2\geq\frac23-0.4397243>0.2269.
\]
\end{corollary}

\begin{proof}
Since $F_1(d)=\sigma(d)$,
\[
 G_2=\sigma(\N)\cup F_2(\N).
\]
Lemma~\ref{lem:sigma-range-zero} implies that $\sigma(\N)$ has density
zero, and therefore $\upperdens(G_2)\leq\theta_2$.  Moreover,
$\Lambda_2=2$, $P_2=4$, $\mathfrak Q(P_2)=6$, and so $\delta_2=1/3$.
Using \eqref{eq:bounded-cofactor-complement},
\[
 \eta_2
 =1-\lowerdens(G_2)-\delta_2
 \geq\frac23-\upperdens(G_2)
 \geq\frac23-0.4397243>0.2269.
\]
\end{proof}

\begin{remark}
This improvement concerns only $A=2$.  The next result gives, for every
$A$, an explicit lower bound for $1-\lowerdens(G_A)$ from the forced
congruences of Section~\ref{subsec:forced-modulus}.
\end{remark}
For $A\geq2$ let $\mathcal K_A$ be the set of moduli $K_{e,d}$ with
$2\leq e\leq A$ and $d\geq1$; it is finite by
Section~\ref{subsec:forced-modulus}.  Put
\[
 V_A=\{N\in\N:k\nmid N\ \text{for every }k\in\mathcal K_A\} .
\]
This is a union of residue classes modulo $W_A=\lcm\mathcal K_A$, so it has
a natural density, and $d(V_A)\geq\varphi(W_A)/W_A>0$.  By inclusion and
exclusion the density is the rational number
\begin{equation}\label{eq:fm-ie}
 d(V_A)=1-\sum_{\varnothing\neq I\subseteq\mathcal K_A}
 \frac{(-1)^{|I|+1}}{\lcm I} .
\end{equation}
\begin{proposition}[Forced-modulus envelope]\label{prop:fm-envelope}
Fix $A\geq2$.  All but $o_A(X)$ of the integers $N\leq X$ in $G_A$ are
divisible by some $k\in\mathcal K_A$.  Consequently
\[
 1-\lowerdens(G_A)\geq d(V_A) .
\]
\end{proposition}
\begin{proof}
Let $N\leq X$ lie in $G_A\cap V_A$.  If $N$ is a value of $F_1$ only, then
$N=\sigma(d)$ lies in $\sigma(\N)$, which has density zero by
Lemma~\ref{lem:sigma-range-zero}.  Otherwise $N=F_e(d)$ with
$2\leq e\leq A$.  Put $m=ed\leq AX$ and $k=K_{e,d}\in\mathcal K_A$.  By
\eqref{eq:fm-congruence}, $N\equiv\sigma(m)\pmod k$, and $k\nmid N$ forces
$k\nmid\sigma(m)$, hence $W_A\nmid\sigma(m)$.  The pair $(e,m)$ determines
$N$, so these $N$ number at most
\[
 A\cdot\#\{m\leq AX:W_A\nmid\sigma(m)\}=o_A(X)
\]
by Lemma~\ref{lem:fixed-modulus-normality}.  Thus all but $o_A(X)$ of the
elements of $V_A\cap[1,X]$ lie outside $G_A$, which gives the density
bound.
\end{proof}
The envelope already accounts for the almost-logarithmic scale in
Theorem~\ref{thm:almost-log-tail}.  Here $\Delta_y$ is as in
Section~\ref{sec:large-values}.
\begin{corollary}\label{cor:fm-envelope-tail}
As $A\to\infty$,
\[
 \Bigl(\frac{e^{-\gamma}}2-o(1)\Bigr)L(A)\leq d(V_A)\leq\Delta_A .
\]
\end{corollary}
\begin{proof}
Every prime $p\leq A$ lies in $\mathcal K_A$: with $e=p$ and $d=1$ one has
$D=\{1\}$, $M=C=g=1$ and $K_{p,1}=p$.  So $V_A$ consists of $A$-rough
integers, and $d(V_A)\leq\Delta_A$.

For the lower bound let $P$ be the parameter of
Lemma~\ref{lem:kernel-tails} with $E=A$.  Then
$\log P=(2+o(1))(\log A)(\log_2A)/\log_3A$, and Mertens' theorem gives
\[
 \Delta_P=\Bigl(\frac{e^{-\gamma}}2+o(1)\Bigr)L(A).
\]
It suffices to show that all but a proportion $o(1)$ of the $P$-rough
integers lie in $V_A$.

Let $N$ be $P$-rough and divisible by $k=K_{e,d}$ for some $2\leq e\leq A$.
Then $k$ is $P$-rough.  Since $e/g$ divides $k$ and $e/g\leq A<P$, we must
have $e/g=1$, that is, $e\mid M$ and $k=C$.  Now $C\geq M\geq e\geq2$, so
$C>P$, and $C\leq(e-1)M$ gives $M>P/A$.

Given $M$, there are at most $\tau(M)$ choices of $e$, since $e\mid M$, and
$e$ and $M$ together determine $D=\{j<e:j\mid M\}$ and hence $C$.  So at
most $\tau(M)$ moduli $k$ share a given $M$, and for each of them the
$P$-rough multiples of $C$ have density $\Delta_P/C\leq\Delta_P/M$.

Since $M\mid\Lambda_A$ and $M>P/A$,
\[
 d\{N:N\ \text{$P$-rough},\ N\notin V_A\}
 \leq\Delta_P\sum_{\substack{M\mid\Lambda_A\\ M>P/A}}\frac{\tau(M)}M
 =\Delta_P\,S(A,P/A)=o(\Delta_P)
\]
by \eqref{eq:fixed-kernel-tail}.  Hence $d(V_A)\geq(1-o(1))\Delta_P$.
\end{proof}

Corollary~\ref{cor:fm-envelope-tail} recovers the lower bound of
Corollary~\ref{cor:fixed-cofactor-defect} with a periodic witness set.

\begin{conjecture}[Weak bounded-cofactor conjecture]
\label{conj:bounded-cofactor-weak}
As $A\to\infty$ through the positive integers,
\[
 \lowerdens(G_A)\longrightarrow1.
\]
Equivalently, $\eta_A=o(1)$.
\end{conjecture}

The equivalence in Conjecture~\ref{conj:bounded-cofactor-weak} follows from
\eqref{eq:bounded-cofactor-complement} and
\eqref{eq:delta-A-asymptotic}.

\begin{remark}\label{rmk:fm-limit}
Corollary~\ref{cor:fm-envelope-tail} shows where the congruence argument
stops.  The set $V_A$ consists of $A$-rough integers, so the forced
congruences say nothing about integers with a prime factor at most $A$
that are missed by $G_A$.  The $A$-rough part of $\N\setminus G_A$ has
density at most $\Delta_A$, which tends to zero.
Conjecture~\ref{conj:bounded-cofactor-weak} is therefore equivalent to the
assertion that the part of $\N\setminus G_A$ with a prime factor at most
$A$ has density tending to zero.  Every element of that part satisfies some
forced congruence and is nevertheless not a value.
\end{remark}

\begin{proposition}[Bounded cofactors and tightness]
\label{prop:tightness-equivalence}
The empirical measures $\nu_X$ defined in \eqref{eq:empirical-measures} form
a tight family on $[0,\infty)$ if and only if
\begin{equation}\label{eq:tightness-criterion}
 \lim_{E\to\infty}\upperdens(\N\setminus G_E)=0.
\end{equation}
Equivalently,
\[
 \lim_{E\to\infty}\lowerdens(G_E)=1.
\]
Thus Conjecture~\ref{conj:bounded-cofactor-weak} is exactly the assertion
that the family $\{\nu_X\}$ is tight.
\end{proposition}

\begin{proof}
The summand $d$ occurs in $F_e(d)$, so $F_e(d)\geq d$.  Hence
\[
 N\in G_E\quad\Longrightarrow\quad f(N)\leq ed\leq EN
 \quad\text{for some }e\leq E.
\]
Using \eqref{eq:exact-representability} to replace the normalization by
$X+O(1)$, we obtain
\[
 \limsup_{X\to\infty}\nu_X((E,\infty))
 \leq\upperdens(\N\setminus G_E).
\]
Thus \eqref{eq:tightness-criterion} implies tightness.

Conversely, fix $T>0$ and an integer $k\geq2$.  Every represented
$N\notin G_E$ with $f(N)\leq TN$ has a witnessing cofactor $e>E$.
Lemma~\ref{lem:moment} therefore gives
\[
 \upperdens(\N\setminus G_E)
 \leq\limsup_{X\to\infty}\nu_X((T,\infty))
      +C_kT^{k+1}E^{1-k}.
\]
If the family is tight, first let $E\to\infty$ and then let $T\to\infty$.
The right-hand side tends to zero.  Finally,
$\upperdens(\N\setminus G_E)=1-\lowerdens(G_E)$ identically.
\end{proof}

Since $r_A(N)>0$ is equivalent, apart from the two unrepresented integers,
to $f(N)\leq AN$, Proposition~\ref{prop:tightness-equivalence} also shows
that the density-one statement \eqref{eq:T} is equivalent to the weak
bounded-cofactor conjecture.

\subsection{Further questions}

The preceding results leave several related problems.

\paragraph{Optimal upper tails.}
Theorem~\ref{thm:almost-log-tail} gives a fixed-level lower bound of order
\[
 \frac{\log_3T}{(\log T)(\log_2T)},
\]
while Theorem~\ref{thm:subexp-growth} reaches the growing threshold
$\exp((1/\sqrt2-o(1))\sqrt{\log_3N\log_4N})$ on
$\gg X/\log_3X$ targets.  No matching upper bounds are known.  It remains to
determine the correct fixed-level tail, the optimal growing scale, and
whether the constant $1/\sqrt2$ reflects a genuine barrier or only the
present divisor-tail estimate.

\paragraph{Bounded-cofactor coverage.}
By Proposition~\ref{prop:tightness-equivalence}, tightness is exactly the
assertion that the finite unions $G_E=\bigcup_{e\leq E}F_e(\N)$ eventually
cover a set of density arbitrarily close to one.  The forced congruences of
Section~\ref{subsec:forced-modulus} account for an $E$-rough part of what
is missed, of density at most $\Delta_E$, and Proposition~\ref{prop:fm-eta}
shows that the coprime part $\delta_E$ is not the main term either.  What
remains is the set of integers with a small prime factor that satisfy every
forced congruence and are still not values.  Nothing in this paper bounds
its density, and tightness is equivalent to that density tending to zero.

\paragraph{Collisions between representations.}
Proposition~\ref{prop:class-first-moment} shows that the number of bounded
representations tends to infinity in every fixed residue class, but many
representation pairs may have the same target.  The missing step is therefore
an upper bound for genuine collisions.  A natural route is to estimate the
second moment of $r_A(N)$, or of a suitably chosen subfamily, on the scale of
the square of its first moment; a sufficiently uniform such estimate would
turn the divergent first moment into density-one coverage.

\paragraph{Limiting distributions.}
Corollary~\ref{cor:limit-law-tails} describes every finite subsequential limit:
it has no atom at zero, has zero in its support, and has unbounded support.
Theorem~\ref{thm:subexp-growth} also forces all positive empirical moments to
diverge.  The remaining questions are whether the family is tight, whether
the full sequence $\nu_X$ converges, whether a subsequential limit is unique,
and whether such a limit has full support on $(0,\infty)$.

\paragraph{Individual cofactor ranges.}
For a fixed $e$ the range $F_e(\N)$ is thin.  The range $F_1(\N)=\sigma(\N)$
has density zero, $\upperdens F_2(\N)\leq0.4398$ by
Proposition~\ref{prop:theta-two}, and $\upperdens F_p(\N)\leq1/p$ for every
prime $p$ by Corollary~\ref{cor:fm-prime-ceiling}.  Coverage of a set of
density close to one must therefore come from the union over $e$, and the
question is how the ranges correlate as $e$ varies.  Whether each
$F_e(\N)$ has positive lower density is known for $e=2$ by Luca and
Pomerance \cite{LucaPomerance} and is open in general.

\begin{thebibliography}{99}

\bibitem{ChenZhao}
Y.-G.~Chen and Q.-Q.~Zhao,
\emph{Nonaliquot numbers},
Publ. Math. Debrecen \textbf{78} (2011), no.~2, 439--442,
\url{https://doi.org/10.5486/PMD.2011.4820}.

\bibitem{MVGoldbach}
H.~L. Montgomery and R.~C. Vaughan,
\emph{The exceptional set in Goldbach's problem},
Acta Arith. \textbf{27} (1975), 353--370,
\url{https://doi.org/10.4064/aa-27-1-353-370}.

\bibitem{jif}
jif,
comment under Erd\H{o}s Problem \#1054 on representability,
April 16, 2026,
\url{https://www.erdosproblems.com/forum/thread/1054}.

\bibitem{jifRepo}
jif,
\emph{Erd\H{o}s Problem 1054 verifier},
GitHub repository,
\url{https://github.com/jif-perso/erdos_1054}.

\bibitem{Kovac}
V.~Kova\v{c},
\emph{An improved bound in Erd\H{o}s Problem \#1054},
May 17, 2026,
\url{https://web.math.pmf.unizg.hr/~vjekovac/files/Erdos_problem_1054.pdf}.

\bibitem{TaoComment}
T.~Tao,
comment under Erd\H{o}s Problem \#1054,
November 1, 2025,
\url{https://www.erdosproblems.com/forum/thread/1054\#post-1636}.

\bibitem{LucaPomerance}
F.~Luca and C.~Pomerance,
\emph{The range of the sum-of-proper-divisors function},
Acta Arith. \textbf{168} (2015), no.~2, 187--199.

\bibitem{Pollack}
P.~Pollack,
\emph{Some arithmetic properties of the sum of proper divisors and the sum
of prime divisors},
Illinois J. Math. \textbf{58} (2014), no.~1, 125--147.

\bibitem{Tenenbaum}
G.~Tenenbaum,
\emph{Introduction to Analytic and Probabilistic Number Theory},
3rd ed., Graduate Studies in Mathematics, vol.~163,
American Mathematical Society, Providence, RI, 2015.

\end{thebibliography}

\end{document}
